\documentclass[11pt]{article}
\usepackage[margin=1in]{geometry}
\usepackage{amsmath,amssymb,amsthm,mathtools}
\usepackage{hyperref}
\usepackage{enumitem}
\hypersetup{colorlinks=true,linkcolor=blue,urlcolor=blue}
\newcommand{\Z}{\mathbb{Z}}
\newcommand{\R}{\mathbb{R}}
\newcommand{\C}{\mathbb{C}}
\newcommand{\im}{\operatorname{im}}
\newcommand{\rank}{\operatorname{rank}}
\newcommand{\GL}{\operatorname{GL}}
\newcommand{\SL}{\operatorname{SL}}
\newcommand{\Sp}{\operatorname{Sp}}
\newcommand{\Tr}{\operatorname{Tr}}
\DeclareMathOperator{\Span}{span}
\DeclareMathOperator{\Hom}{Hom}
\newenvironment{solution}[1]{\medskip\noindent\textbf{P#1.}\ }{\par\medskip}
\title{The Cobordism Functor, the Theta Register, and Root Systems\\[4pt]
\large Solutions}
\author{}
\date{}
\begin{document}
\maketitle

Solutions to the companion \emph{Problems}, in its numbering; theorem and
section numbers refer to \path{docs/design/cobordism_functor_reading.tex}.
Where a problem asks for a value the verifier reports, the recorded value is
quoted.

\part*{The primer}

\section*{Cochains, the Laplacian and the harmonic space (Section~2)}

\begin{solution}{1}
On a simplex $\sigma = [v_0\dots v_{p+q+1}]$, with $(\alpha\cup\beta)[v_0\dots v_{p+q}]
= \alpha[v_0\dots v_p]\,\beta[v_p\dots v_{p+q}]$,
\[
  d(\alpha\cup\beta)(\sigma) = \sum_{i=0}^{p+q+1}(-1)^i\,(\alpha\cup\beta)[v_0\dots\hat v_i\dots v_{p+q+1}].
\]
Split at $i\le p$ and $i\ge p+1$. The first group is $(d\alpha\cup\beta)(\sigma)$
minus its $i=p+1$ term; the second is $(-1)^p(\alpha\cup d\beta)(\sigma)$ minus
its $j=0$ term. Those two omitted terms are $(-1)^{p+1}\alpha[v_0..v_p]\beta[v_{p+1}..]$
and $(-1)^p\alpha[v_0..v_p]\beta[v_{p+1}..]$, which cancel. Hence
$d(\alpha\cup\beta) = d\alpha\cup\beta + (-1)^p\alpha\cup d\beta$. For closed
$\alpha,\beta$: $\langle\alpha|_{\partial W}\cup\beta|_{\partial W},[\partial W]\rangle
= \langle\alpha\cup\beta,\partial[W]\rangle = \langle d(\alpha\cup\beta),[W]\rangle = 0$.
\end{solution}

\begin{solution}{2}
Every edge starts and ends at $v$, so $\partial_1 = 0$ and $d_0 = \partial_1^{\mathsf T} = 0$.
With the ordering $(a,b,c)$ and $(U,L)$, $\partial_2 = \begin{psmallmatrix}1&-1\\1&-1\\-1&1\end{psmallmatrix}$
and $d_1 = \partial_2^{\mathsf T} = \begin{psmallmatrix}1&1&-1\\-1&-1&1\end{psmallmatrix}$;
$d_1d_0 = 0$ trivially. $\ker d_1 = \{\alpha_c = \alpha_a+\alpha_b\}$ has
dimension $2$ and $\im d_0 = 0$, so $b_1 = 2$; $b_0 = 1$ ($C^0 = \C$, $d_0 = 0$);
$b_2 = 2-\rank d_1 = 1$. With $M_0 = 1$, $\delta_1 = M_0^{-1}d_0^{\mathsf T}M_1 = 0$,
so $L_1 = \delta_2d_1 = M_1^{-1}d_1^{\mathsf T}M_2d_1$. For $\alpha$ with
$x = \alpha_a+\alpha_b-\alpha_c$, $d_1\alpha = (x,-x)$ and
$d_1^{\mathsf T}M_2(x,-x)^{\mathsf T} = (m_U+m_L)\,x\,(1,1,-1)^{\mathsf T}$. If
$m_L = -m_U$ this vanishes for every $\alpha$: $L_1 = 0$ and $\ker L_1 = C^1$,
dimension $3$. The closedness certificate fails: $d_1\alpha = (x,-x)\ne0$ for
$\alpha\notin\ker d_1$ although $L_1\alpha = 0$; coclosedness holds vacuously
and the span test would report dimension $3$ against $b_1 = 2$. Theorem~2.2
needs positive Hermitian masses, for which $\langle h,L_1h\rangle_{M_1} = \|d_1h\|^2_{M_2}+\|\delta_1h\|^2_{M_0}$
forces $d_1h = 0$; with $M_2$ indefinite the ``energy'' $(m_U+m_L)|x|^2$ can
vanish on a nonzero $d_1h$. This is the Lorentzian caveat after Theorem~2.2.
\end{solution}

\begin{solution}{3}
Two representatives $\omega,\omega'$ of the same class $\alpha$ differ by
$d\eta$; by Lemma~2.1, $\langle d\eta,z\rangle = \langle\eta,\partial z\rangle = 0$
for every cycle $z$, so the periods agree. J1 (fixed $M$) is the corollary to
Theorem~4.2: $M$ is built from periods only. J2 (the representative moves by
an $O(1)$ amount) measures that the harmonic condition $\delta\omega = 0$ depends
on the masses, so the chosen cochain in the class changes. J3 (the difference
is exactly in $\im d_0$) measures the proof of Theorem~2.4: if the difference
had a non-exact part, the two cochains would lie in different classes, which
cannot happen for a change of lengths on a fixed complex; it would signal a
certificate failure (the band was not the harmonic space) or a singular $P_A$.
\end{solution}

\section*{Restriction and the monodromy (Sections~3--4)}

\begin{solution}{4}
Lefschetz gives $H^0(W,\partial W)\cong H_3(T^2)=0$, $H^1(W,\partial W)\cong H_2(T^2)=\Z$,
$H^2(W,\partial W)\cong H_1(T^2)=\Z^2$, $H^3(W,\partial W)\cong\Z$. With $H^0(W)=\Z$,
$H^0(\partial W)=\Z^2$, $H^1(W)=\Z^2$, $H^1(\partial W)=\Z^4$, $H^2(W)=\Z$, $H^2(\partial W)=\Z^2$
the sequence is
\[
  0\to 0\to\Z\to\Z^2\to\Z\to\Z^2\xrightarrow{\ \iota^*\ }\Z^4\to\Z^2\to\Z\to\Z^2\to\Z\to 0 .
\]
$H^0(W)\to H^0(\partial W)$ is the diagonal, rank $1$, so its cokernel $\Z$ maps
\emph{onto} $H^1(W,\partial W)=\Z$, so $H^1(W,\partial W)\to H^1(W)$ is zero, so
$\iota^*$ is injective: $\rank\iota^* = 2 = \tfrac12\cdot 4$. Consistency: the
next map has rank $4-2=2$, onto $\Z^2$, so $H^2(W)\to H^2(\partial W)$ is injective
of rank $1$ and $H^2(\partial W)\to H^3(W,\partial W)$ has rank $1$, onto.
\end{solution}

\begin{solution}{5}
$W$ retracts onto $T^2$ and both ends restrict by the identity, so
$L_W = \{(p,p): p\in\C^2\}$, the diagonal. For $\Omega_\partial = J\oplus(-J)$:
$(p,p)^{\mathsf T}\Omega_\partial(q,q) = p^{\mathsf T}Jq - p^{\mathsf T}Jq = 0$. Dimension $2$
in $\C^4$: Lagrangian.
\end{solution}

\begin{solution}{6}
$J\sigma = \begin{psmallmatrix}0&1\\-1&0\end{psmallmatrix}\begin{psmallmatrix}0&1\\1&0\end{psmallmatrix}
= \begin{psmallmatrix}1&0\\0&-1\end{psmallmatrix}$, so
$\sigma^{\mathsf T}J\sigma = \begin{psmallmatrix}0&1\\1&0\end{psmallmatrix}\begin{psmallmatrix}1&0\\0&-1\end{psmallmatrix}
= \begin{psmallmatrix}0&-1\\1&0\end{psmallmatrix} = -J$. Hence
$(\sigma\oplus\sigma)^{\mathsf T}(J\oplus J)(\sigma\oplus\sigma) = -(J\oplus J)$ while
$\det(\sigma\oplus\sigma) = (-1)^2 = +1$. The $2\times2$ identity would predict
$+(J\oplus J)$. Each block reverses the form independently; the reversal is a
property of each torus's relabelling, so the sign must be declared per torus.
(The code first tied it to $\det M$ and this counterexample corrected it.)
\end{solution}

\begin{solution}{7}
The far marking is $(\varphi A,\varphi B)$, so $P_B = \varphi^*P_A$ and $M = \varphi^*$.
Swap: exchanges $A,B$; $\varphi^* = \begin{psmallmatrix}0&1\\1&0\end{psmallmatrix}$,
$\det=-1$. Dehn twist $(x,y)\mapsto(x+y,y)$: $A\mapsto A$, $B\mapsto A+B$, so
$\varphi^* = \begin{psmallmatrix}1&1\\0&1\end{psmallmatrix}$, $\det=+1$. On the grid
with diagonal $(i,j)\!-\!(i{+}1,j{+}1)$ the swap sends that diagonal to
$(j,i)\!-\!(j{+}1,i{+}1)$, a diagonal of the same family: simplicial. The twist
$(i,j)\mapsto(i{+}j,j)$ sends the diagonal to $(i{+}j,j)\!-\!(i{+}j{+}2,j{+}1)$, a
displacement $(2,1)$, which is not an edge: not simplicial. Automorphisms of a
finite triangulation have finite order, and a twist has infinite order, so no
automorphism realizes it. A flip pass changes the triangulation: flipping every
edge of one family of the far torus produces the same grid drawn in a sheared
basis, i.e.\ the image of the grid under the lattice map $L$. Each flip is a
$2$--$2$ Pachner move, which is the boundary of one tetrahedron; stacking one
tetrahedron per flipped edge gives a layer whose near face is the old grid and
whose far face is the sheared grid. The far marking is read through $L$, so the
collar's relation is $M = L^{\mathsf T}$; the measured passes give
$T_B = \begin{psmallmatrix}1&1\\0&1\end{psmallmatrix}$, $T_A$ and $S$ (F2).
\end{solution}

\begin{solution}{8}
$1.6\times10^{-14}$ measures the corollary to Theorem~4.2 (metric independence
of $M$), whose mechanism is Theorem~2.4: periods of a fixed class do not move.
The $0.7$ move with an exact difference measures Theorem~2.4's proof: the two
harmonic representatives of one class differ by a coboundary. A nonzero
non-exact part would mean the representatives lay in \emph{different} classes
--- impossible for a change of lengths on a fixed complex --- so it would
signal that the band was not the harmonic space (a certificate failure) or
that $P_A$ was singular and the canonical normalization failed.
\end{solution}

\begin{solution}{9}
$M = P_BP_A^{-1} = \operatorname{diag}(2,\tfrac12)$, $\det M = 1$, and for $2\times2$
matrices $\det = 1$ is symplecticity, but $\tfrac12\notin\Z$. Theorem~4.2 assumes
that the two restrictions identify the bulk lattice with each boundary lattice
$H^1(M_i;\Z)$; then $M$ is the matrix of a $\Z$-module isomorphism, hence
integral with unit determinant. Invertibility over $\C$ says only that the
restriction is injective on a basis, not that it hits the integral basis.
Connected sum: by Mayer--Vietoris with $W_1\cap W_2\simeq S^2$ (Proposition~16.1),
$H^1(W) = H^1(T^2\times I)\oplus H^1(S^1\times D^2) = \Z^2\oplus\Z$, and the new
class is supported away from the boundary tori, so its periods on both ends
vanish. With a harmonic basis of three forms, $P_A$ and $P_B$ are $2\times3$ with
the invisible class in $\ker P_A\cap\ker P_B$. The identifiability proposition:
$B = TA$ has a unique solution on $\operatorname{im}A$ exactly when $A$ is onto
and $\ker A\subseteq\ker B$; both hold, so $M$ is the unique map with
$P_B = MP_A$, and it equals the product collar's $I$.
\end{solution}

\section*{Choi vectors and composition (Sections~5--6)}

\begin{solution}{10}
Let $\pi_A:L\to A$ be the projection. $\ker\pi_A = L\cap(0\oplus B)$, so $\pi_A$
is injective iff that is zero, and it is onto by hypothesis; then $\pi_A|_L$ is
an isomorphism and $T := \pi_B\circ(\pi_A|_L)^{-1}$ is the unique linear map
with $L = \{(x,Tx)\}$. Conversely a graph meets $0\oplus B$ trivially and
projects onto $A$. $\dim L = \dim A$.
\end{solution}

\begin{solution}{11}
$|I\rangle\!\rangle = e_1\otimes e_1^* + e_2\otimes e_2^*$, the Bell vector, Schmidt rank
two; $|\sigma\rangle\!\rangle = e_2\otimes e_1^* + e_1\otimes e_2^*$, also maximally
entangled. The product collar's Choi vector is maximally entangled because the
identity channel correlates input and output perfectly: under Choi, unitary
channels correspond exactly to maximally entangled operator vectors. This is
operator-vector information (Section~29), not entanglement of a physical
two-torus state.
\end{solution}

\begin{solution}{12}
With $U = W_1$, $V = W_2$, $U\cap V = T^2\times\{1\}$: $H^0(U)\oplus H^0(V)\to H^0(T^2)$
is $(c_1,c_2)\mapsto c_1-c_2$, onto, so $H^1(W)\to H^1(W_1)\oplus H^1(W_2)$ is
injective with image the kernel of $(\alpha_1,\alpha_2)\mapsto\alpha_1|_{T^2}-\alpha_2|_{T^2}$,
i.e.\ compatible pairs $\cong H^1(T^2) = \Z^2$. The far ends read the compatible
pair's two restrictions, both the identity: $L_W$ is the diagonal and
$M = I\cdot I = I$. The result is the product collar $T^2\times[0,2]$.
\end{solution}

\begin{solution}{13}
($\subseteq$) A class $\alpha$ on $W$ restricts to $\alpha_1,\alpha_2$ on $W_1,W_2$
which agree on $M_1$ (both are $\alpha|_{M_1}$); then $(\alpha|_{M_0},\alpha|_{M_1})\in L_{W_1}$
and $(\alpha|_{M_1},\alpha|_{M_2})\in L_{W_2}$, so $(\alpha|_{M_0},\alpha|_{M_2})\in L_{W_2}\circ L_{W_1}$
with $x_1 = \alpha|_{M_1}$. ($\supseteq$) Given $(x_0,x_1)\in L_{W_1}$ and $(x_1,x_2)\in L_{W_2}$,
choose $\alpha_1\in H^1(W_1)$ restricting to $(x_0,x_1)$ and $\alpha_2\in H^1(W_2)$
restricting to $(x_1,x_2)$. They agree on $M_1$, so $(\alpha_1,\alpha_2)$ lies in
the kernel of $H^1(W_1)\oplus H^1(W_2)\to H^1(M_1)$, which by exactness is the
image of $H^1(W)$: some $\alpha$ restricts to both, hence to $(x_0,x_2)$. Exactness
at $H^1(W_1)\oplus H^1(W_2)$ is the step used. The connecting map
$H^0(M_1)\to H^1(W)$ makes $\alpha$ non-unique (two extensions differ by an
invisible class), but every extension has the same restrictions, so the
relation is unaffected; only the harmonic space, not $L_W$, sees the extra
class. On a simplicial complex the sequence requires one cell structure on
the seam, which is why separately seeded collars do not compose by this
theorem.
\end{solution}

\begin{solution}{14}
The $2L$-layer collar is seeded as one complex; nothing identifies interface
cells of two $L$-layer collars, and no harmonic form is required to be
continuous across a seam, so Mayer--Vietoris is never exercised. The check
verifies $M=\varphi^*$ on two hosts and the relation $\sigma^2 = I$ in the
group. The automorphisms realizable on the square grid, $\{\pm I,\pm\sigma\}\cong\Z/2\times\Z/2$,
commute, so among automorphisms no experiment distinguishes $M(W_2\circ W_1)$
from $M(W_1\circ W_2)$. The flip passes stack two layers on one host, so the
far marking is read through $L_2L_1$ and the relation composes on a single
complex: $T_AT_B = \begin{psmallmatrix}1&1\\1&2\end{psmallmatrix}$ and
$T_BT_A = \begin{psmallmatrix}2&1\\1&1\end{psmallmatrix}$ differ (F4), and their
Weil matrices are at projective distance $1$ (T9). This is Theorem~6.1 on one
complex with the seam being a face layer of the same triangulation. It still
does not glue two separately seeded cobordisms along a shared boundary torus:
that requires identifying interface cells and frames of two complexes.
\end{solution}

\section*{What the metric controls (Section~7)}

\begin{solution}{15}
Under the gauge the dressed metric transforms by $G\mapsto\rho G\rho^{-1}$ and
the primal frame by $\Phi\mapsto\rho\Phi$. Then $\Phi^{\mathsf T}G\Phi\mapsto
\Phi^{\mathsf T}\rho^{\mathsf T}\rho G\rho^{-1}\rho\Phi = \Phi^{\mathsf T}\rho^2G\Phi$
($\rho$ diagonal), which differs from $\Phi^{\mathsf T}G\Phi$ unless $\rho^2=I$.
The dual frame transforms by $\tilde\Phi\mapsto\rho^{-1}\tilde\Phi$, so
$\tilde\Phi^{\mathsf T}G\Phi\mapsto\tilde\Phi^{\mathsf T}\rho^{-1}\rho G\rho^{-1}\rho\Phi
= \tilde\Phi^{\mathsf T}G\Phi$. At zero phases the dual connection equals the
connection, $\Phi^\vee = \Phi$, and $\tilde\Phi = G\Phi^\vee B^{-\mathsf T} = ZB^{-\mathsf T}$
with $B = \Phi^{\mathsf T}Z$; then $\tilde\Phi^{\mathsf T}MZ = B^{-1}Z^{\mathsf T}MZ = B^{-1}B = I$
and the two Grams coincide. Measured: $0.76$ against $1.3\times10^{-15}$;
agreement at zero phases to $3.6\times10^{-16}$.
\end{solution}

\begin{solution}{16}
Every symbol in $D$ is boundary data: $G_A,G_B$ are the own-boundary Grams
(functions of the boundary lengths), the markings are fixed, and $M,M^\vee$
are period transports, which by Theorem~2.4 depend only on the classes. No
argument depends on interior lengths, so no bulk relaxation changes $D$; a
verifier that ``expects the defect to fall'' would be measuring nothing
(this is why V2 keeps it as a control). $D = 0$ says the bilinear pairing
$G$ is preserved, $(M^\vee)^{\mathsf T}G_BM = G_A$, with $G$ a complex bilinear
form; unitarity is preservation of a positive Hermitian form,
$T^\dagger Q_{\rm out}T = Q_{\rm in}$ (Section~29). A complex bilinear form is
not an inner product, so the two statements are different tests.
\end{solution}

\section*{Weights, phase space and quantization (Sections~8--9)}

\begin{solution}{17}
A character $\chi:\mathfrak h/\Lambda\to U(1)$ lifts to a homomorphism
$\mathfrak h\to\R$, $h\mapsto\langle\mu,h\rangle$, trivial on $\Lambda$ iff
$\langle\mu,\Lambda\rangle\subseteq\Z$, i.e.\ $\mu\in\Lambda^*$. On $V\otimes V'$,
$T$ acts diagonally, $t\cdot(v\otimes v') = \chi_\mu(t)\chi_{\mu'}(t)\,v\otimes v'
= \chi_{\mu+\mu'}(t)\,v\otimes v'$: weights add. For $T_1\times T_2$ a character
is a pair of characters, so $\Lambda^* = \Lambda_1^*\oplus\Lambda_2^*$. The
space spanned by all characters of $J(\Sigma) = H^1(\Sigma;\R)/H^1(\Sigma;\Z)$
is $L^2$ of a torus, infinite-dimensional, with no distinguished
finite-dimensional subspace. Theta quantization chooses a complex structure
(the period matrix $\Omega$), a positive line bundle $L$ whose curvature is the
intersection form (the principal polarization) and a level $k$; the
holomorphic sections $H^0(J,L^k)$ are finite-dimensional ($k^g$) because
holomorphy plus the quasi-periodicity fixes every Fourier coefficient from
$k^g$ residue classes (Theorem~11.2). A disjoint union gives a product
Jacobian with the external product bundle, whose sections are the tensor
product (Proposition~11.1).
\end{solution}

\begin{solution}{18}
A gauge transformation adds $d\chi$ to $\phi$, and $\oint_\gamma d\chi = 0$
around any cycle (Lemma~2.1), so $a,b$ are unchanged. A large gauge
transformation is a closed cochain with periods in $2\pi\Z$ (the exponential
of $i$ times it is a well-defined phase), so it shifts $(a,b)$ by integers.
Hence the physical configuration space is $H^1(\Sigma;\R)$ modulo
$H^1(\Sigma;\Z)$, a torus of dimension $b_1 = 2g$ with coordinates
$a_i,b_i\in\R/\Z$. On the basis dual to a symplectic marking the intersection
form is $\begin{psmallmatrix}0&I\\-I&0\end{psmallmatrix}$ and the $2$-form
$\sum_ida_i\wedge db_i$ takes the value $1$ on each $(a_i,b_i)$ square and $0$
on the others: its periods on $H_2$ of the torus are the entries of the
intersection form, a unimodular integral class, the principal polarization.
Shifting $b\mapsto b+e_i$ gives $z\mapsto z+e_i$; shifting $a\mapsto a+e_i$
gives $z\mapsto z-\Omega e_i$. Together these generate $\Z^g+\Omega\Z^g$, and
the map $(a,b)\mapsto z = b-\Omega a$ is a real-linear isomorphism
$\R^{2g}\to\C^g$ because $\operatorname{Im}\Omega$ is invertible. So $z$
descends to an isomorphism of $J(\Sigma)$ with the complex torus
$\C^g/(\Z^g+\Omega\Z^g)$; the complex structure is the choice of $\Omega$.
\end{solution}

\begin{solution}{19}
$d_1d_0 = -(v-1)(u-1)+(u-1)(v-1) = 0$. If $(u,v)\ne(1,1)$ then $d_0\ne0$ has
rank $1$, so $H^0 = \ker d_0 = 0$; $d_1\ne0$ has rank $1$, so $\ker d_1$ has
dimension $1 = \rank d_0$ and $H^1 = \ker d_1/\im d_0 = 0$; $H^2 = \C/\im d_1 = 0$.
The Euler characteristic $1-2+1 = 0$ is consistent. So inserting the holonomy
phases into the differential (a twisted pencil) empties the register: there is
no harmonic space left to carry a state. The phases are configuration data
(the point $z$ of the Jacobian), read as Wilson loops of a configuration on
the neutral complex; the Hodge assembly must stay untwisted.
\end{solution}

\section*{Theta functions and the tensor register (Sections~10--11)}

\begin{solution}{20}
$z\mapsto z+1$ multiplies the $n$-th term by $e^{2\pi i k(n+j/k)} = e^{2\pi i(kn+j)} = 1$.
$z\mapsto z+\tau$: the exponent gains $2\pi ik(n+j/k)\tau$; write
$k(n+j/k)^2\tau + 2k(n+j/k)\tau = k[(n+1+j/k)^2-1]\tau$ and
$2\pi ik(n+j/k)z = 2\pi ik(n+1+j/k)z - 2\pi ikz$; reindexing $n+1\to n$ gives
$\theta_j(z+\tau) = e^{-\pi ik\tau-2\pi ikz}\theta_j(z)$.
\end{solution}

\begin{solution}{21}
Under $z\mapsto z+1$ both factors are unchanged. Under $z\mapsto z+\tau$:
$|\theta_j|$ is multiplied by $|e^{-\pi ik\tau-2\pi ikz}| = e^{\pi k\operatorname{Im}\tau+2\pi k\operatorname{Im}z}$,
while $e^{-\pi k(\operatorname{Im}z+\operatorname{Im}\tau)^2/\operatorname{Im}\tau}
= e^{-\pi k(\operatorname{Im}z)^2/\operatorname{Im}\tau}\,e^{-2\pi k\operatorname{Im}z}\,e^{-\pi k\operatorname{Im}\tau}$.
The product of the two extra factors is $1$.
\end{solution}

\begin{solution}{22}
Shift the parallelogram so no zero lies on its boundary. The winding number
of $\theta_j$ around it is $\frac{1}{2\pi i}\oint d\log\theta_j$. The two
vertical sides cancel by periodicity. On the top side $d\log\theta_j(z+\tau)
= d\log\theta_j(z) - 2\pi ik\,dz$, traversed backwards, so the horizontal
sides contribute $\int_0^1 2\pi ik\,dz = 2\pi ik$: exactly $k$ zeros. For $k=2$:
pairing $n\leftrightarrow -n-1$ in the sum shows $\theta_1(\tfrac14) = \theta_1(\tfrac34)=0$
(the paired terms carry $e^{\pi i(n+\frac12)}$ and its negative), and
$\theta_0(\tfrac14+\tfrac\tau2) = \theta_0(\tfrac34+\tfrac\tau2) = 0$ (there the paired
terms carry $(-1)^n$ and $(-1)^{-n-1}$).
\end{solution}

\begin{solution}{23}
$\theta_j(z+\tfrac1k)$: the term gains $e^{2\pi i(n+j/k)} = e^{2\pi ij/k}$, independent
of $n$. $\theta_j(z+\tfrac\tau k)$: the exponent gains $2\pi i(n+j/k)\tau$; with
$k(n+j/k)^2 + 2(n+j/k) = k(n+\tfrac{j+1}{k})^2 - \tfrac1k$ and
$2\pi ik(n+j/k)z = 2\pi ik(n+\tfrac{j+1}k)z - 2\pi iz$ this is
$e^{-\pi i\tau/k-2\pi iz}\theta_{j+1}(z)$. Then $\mathsf Z\mathsf X\theta_j
= e^{2\pi i(j+1)/k}(\cdots)\theta_{j+1}$ and $\mathsf X\mathsf Z\theta_j = e^{2\pi ij/k}(\cdots)\theta_{j+1}$
with the same prefactor, so $\mathsf Z\mathsf X = e^{2\pi i/k}\mathsf X\mathsf Z$.
\end{solution}

\begin{solution}{24}
For $\Omega = \operatorname{diag}(\tau_1,\tau_2)$ the quadratic form is
$\tau_1 m_1^2+\tau_2 m_2^2$ and the linear term $m_1z_1+m_2z_2$, so
$\exp(\cdots)$ is a product of a function of $(n_1,z_1)$ and one of $(n_2,z_2)$
and the double sum factorizes. An off-diagonal $\Omega_{12}$ contributes
$2\pi ik\,m_1m_2\Omega_{12}$, which couples the indices.
\end{solution}

\begin{solution}{25}
Write $\theta_j(x+\tau y) = \sum_n e^{\pi ik(n+j/k)^2\tau+2\pi ik(n+j/k)(x+\tau y)}$.
The $x$-frequencies $k(n+j/k) = kn+j$ are distinct integers for distinct
$(n,j)$, so $\int_0^1\overline{\theta_j}\theta_{j'}dx$ vanishes unless $j = j'$
and pairs equal $n$: it equals $\sum_n|e^{\pi ik(n+j/k)^2\tau+2\pi ik(n+j/k)\tau y}|^2
= \sum_n e^{-2\pi k\operatorname{Im}\tau[(n+j/k)^2+2(n+j/k)y]}$. Multiplying by
$e^{-2\pi k y^2\operatorname{Im}\tau}$ completes the square:
$\sum_ne^{-2\pi k\operatorname{Im}\tau\,(n+j/k+y)^2}$. Integrating $y$ over $[0,1]$
and summing over $n$ tiles the real line once, giving
$\int_\R e^{-2\pi k\operatorname{Im}\tau\,t^2}dt = (2k\operatorname{Im}\tau)^{-1/2}$,
cancelled by the prefactor. The weight is the square of
$e^{-\pi k(\operatorname{Im}z)^2/\operatorname{Im}\tau}$ because
$\operatorname{Im}z = y\operatorname{Im}\tau$. For general $g$, with
$z = x+\Omega y$ and $Y = \operatorname{Im}\Omega$, the same steps give
$\int_{\R^g}e^{-2\pi k t^{\mathsf T}Yt}dt = (2k)^{-g/2}(\det Y)^{-1/2}$, hence
the prefactor $(2k)^{g/2}\sqrt{\det Y}$.
\end{solution}

\section*{The geometric operator algebra (Section~12)}

\begin{solution}{26}
$Z_rX_s|j\rangle = Z_r|j+\epsilon_s\rangle = \omega^{j_r+\delta_{rs}}|j+\epsilon_s\rangle$
and $X_sZ_r|j\rangle = \omega^{j_r}|j+\epsilon_s\rangle$, so $Z_rX_s = \omega^{\delta_{rs}}X_sZ_r$.
Iterating, $Z^qX^{p'} = \omega^{q\cdot p'}X^{p'}Z^q$. Then
$D_{p,q}D_{p',q'} = X^pZ^qX^{p'}Z^{q'} = \omega^{q\cdot p'}X^{p+p'}Z^{q+q'}$.
For the adjoint, $X^\dagger = X^{-1}$ and $Z^\dagger = Z^{-1}$, so
$D_{p,q}^\dagger = Z^{-q}X^{-p} = \omega^{(-q)\cdot(-p)}X^{-p}Z^{-q} = \omega^{p\cdot q}D_{-p,-q}$.
Associativity: $(D_{p,q}D_{p',q'})D_{p'',q''} = \omega^{q\cdot p'+(q+q')\cdot p''}D_{p+p'+p'',\dots}$
and $D_{p,q}(D_{p',q'}D_{p'',q''}) = \omega^{q'\cdot p''+q\cdot(p'+p'')}D_{\dots}$;
both exponents equal $q\cdot p'+q\cdot p''+q'\cdot p''$, the cocycle
condition for the bilinear form $(p,q),(p',q')\mapsto q\cdot p'$.
\end{solution}

\begin{solution}{27}
$D_{p,q}^\dagger D_{p',q'} = \omega^{p\cdot q}D_{-p,-q}D_{p',q'} = \omega^{p\cdot q-q\cdot p'}D_{p'-p,q'-q}$.
If $p'\ne p$ the operator $X^{p'-p}Z^{q'-q}$ maps $|j\rangle$ to a multiple of
$|j+p'-p\rangle\ne|j\rangle$: zero diagonal, zero trace. If $p' = p$ and
$q'\ne q$ the trace is $\sum_j\omega^{(q'-q)\cdot j}$, a sum of a nontrivial
character of $(\Z/k)^g$, which vanishes. If both agree the trace is $d$. Hence
the $D_{p,q}$ are orthogonal for $\langle A,B\rangle = \Tr(A^\dagger B)$; there
are $d^2 = \dim\operatorname{End}(\mathcal Q_k)$ of them, so they are a basis
and the coefficient of $T$ on $D_{p,q}$ is $\Tr(D_{p,q}^\dagger T)/d$. For
$a = \sum a_{p,q}e_{p,q}$, $b = \sum b_{p',q'}e_{p',q'}$, the product law gives
$ab = \sum_{p,q,p',q'}a_{p,q}b_{p',q'}\omega^{q\cdot p'}e_{p+p',q+q'}$; collecting
$r = p+p'$, $s = q+q'$ gives the stated formula. Example ($k = 2$, $\omega = -1$):
$a\star b$ has the single term $r = (1), s = (1)$ from $p = 1,q = 0,p' = 0,q' = 1$
with phase $\omega^{0\cdot0} = 1$: $a\star b = e_{1,1} = XZ$. $b\star a$: $p = 0,q = 1,p' = 1,q' = 0$
with phase $\omega^{1\cdot1} = -1$: $b\star a = -e_{1,1} = -XZ$. Indeed
$ZX = \omega XZ = -XZ$.
\end{solution}

\begin{solution}{28}
$|0\rangle\langle0| = \tfrac12(I+Z)$, so $a = \tfrac12e_{0,0}+\tfrac12e_{0,1}$ and
$\operatorname{tr}_G(a) = 2\cdot\tfrac12 = 1$. $|+\rangle\langle+| = \tfrac12(I+X)$, so
$e = \tfrac12e_{0,0}+\tfrac12e_{1,0}$. By the twisted convolution,
$a\star e = \tfrac14e_{0,0}+\tfrac14e_{1,0}+\tfrac14e_{0,1}+\tfrac14\omega^{1\cdot1}e_{1,1}$,
whose $e_{0,0}$ coefficient is $\tfrac14$; $\operatorname{tr}_G(a\star e) = 2\cdot\tfrac14 = \tfrac12$,
the probability of the outcome $+$ on $|0\rangle$. Positivity of $a$: $a = b^*b$
with $b = a$ since $a\star a = a$ (a projector; check: the $e_{0,0}$ and $e_{0,1}$
coefficients of $a\star a$ are $\tfrac14+\tfrac14 = \tfrac12$ each, using
$e_{0,1}\star e_{0,1} = e_{0,0}$).
\end{solution}

\begin{solution}{29}
In the basis $|j_1\rangle\otimes|j_2\rangle$ of Proposition~11.1, $X_r$ and $Z_r$ act
on the factor containing $r$ and trivially on the other, so
$X^{(p_1,p_2)}Z^{(q_1,q_2)} = (X^{p_1}Z^{q_1})\otimes(X^{p_2}Z^{q_2})$. The
vectorization inner product is $\langle\!\langle A|B\rangle\!\rangle = \Tr(A^\dagger B)$,
so P27 gives orthonormality of $|D_{p,q}\rangle\!\rangle/\sqrt d$.
With $|T\rangle\!\rangle = \sum_{ij}T_{ij}|i\rangle_{\rm out}\otimes|j\rangle_{\rm ref}$,
$|T_2\rangle\!\rangle\otimes|T_1\rangle\!\rangle = \sum T_{2,ij}T_{1,kl}|i\rangle|j\rangle|k\rangle|l\rangle$;
contracting the reference factor of $T_2$ against the output factor of $T_1$
(pairing $\langle j|k\rangle$) leaves $\sum_{il}(\sum_jT_{2,ij}T_{1,jl})|i\rangle|l\rangle = |T_2T_1\rangle\!\rangle$.
For matrix units, $E_{ij}E_{kl} = \delta_{jk}E_{il}$ is exactly the contraction.
The identity holds for unnormalized operator vectors: $\|T_2T_1\|_F$ is not
$\|T_2\|_F\|T_1\|_F$ in general, so renormalizing each factor to a unit state
and then contracting would produce a vector whose norm no longer records the
success amplitude of the composite (for a projector followed by a projector,
the surviving weight).
\end{solution}

\section*{The modular action and the Weil representation (Section~13)}

\begin{solution}{30}
$\pi ik(n+j/k)^2(\tau+1)$ adds $\pi i(kn^2 + 2nj + j^2/k)$; $e^{2\pi inj}=1$
always, $e^{\pi ikn^2} = (-1)^{kn^2} = 1$ for $k$ even, leaving $e^{\pi ij^2/k}$.
For $k$ odd, $e^{\pi ikn^2} = (-1)^n$ depends on $n$, so $\theta_j(z;\tau+1)$ is
not a multiple of $\theta_j(z;\tau)$ (numerically the ratios $|\theta_j(\tau+1)/\theta_j(\tau)|$
are not even constant in $j$); it is a theta function with characteristic
shifted by $\tfrac12$, and $\Theta_k$ is invariant only under a finite-index
subgroup (the theta group generated by $S$ and $T^2$).
\end{solution}

\begin{solution}{31}
Write $\theta_j(z/\tau;-1/\tau) = \sum_n f(n+j/k)$ with
$f(x) = e^{-\pi ikx^2/\tau + 2\pi ikxz/\tau}$. Poisson: $\sum_n f(n+a) = \sum_m\hat f(m)e^{2\pi ima}$.
The Gaussian integral $\int e^{-\pi ax^2+2\pi ibx}dx = a^{-1/2}e^{-\pi b^2/a}$ with
$a = ik/\tau$, $b = kz/\tau - m$ gives
$\hat f(m) = (\tau/ik)^{1/2}e^{\pi ikz^2/\tau}e^{-2\pi izm}e^{\pi im^2\tau/k}$.
Write $m = kn'+l$, $l\in\Z/k$: $m^2/k = k(n'+l/k)^2$ and $zm = k(n'+l/k)z$, so
\[
  \theta_j(z/\tau;-1/\tau) = (\tau/ik)^{1/2}e^{\pi ikz^2/\tau}\sum_l e^{2\pi ijl/k}\theta_l(-z;\tau)
  = (-i\tau)^{1/2}e^{\pi ikz^2/\tau}\sum_l \frac{e^{-2\pi ijl/k}}{\sqrt k}\theta_l(z;\tau),
\]
using $\theta_l(-z) = \theta_{-l}(z)$ and $(\tau/ik)^{1/2} = (-i\tau)^{1/2}k^{-1/2}$
(principal branches, $\operatorname{Im}\tau>0$). So $\rho_k(S)_{jl} = k^{-1/2}e^{-2\pi ijl/k}$
exactly, with automorphy factor $(-i\tau)^{1/2}$ and no extra constant; at $k=2$,
$e^{-\pi ijl} = (-1)^{jl}$ gives the Hadamard matrix, which is real, so the
conjugate sign is invisible there. At $k = 4$ the entries $e^{\mp\pi ijl/2}$
differ, and a numerical evaluation of both sides at one $(z,\tau)$ shows the
residual $3\times10^{-16}$ for $e^{-2\pi ijl/k}$ against $1.5$ for the
conjugate: the sign is forced.
\end{solution}

\begin{solution}{32}
$(\rho(S)^2)_{jm} = \frac1k\sum_l e^{-2\pi i(j+m)l/k} = \delta_{j,-m}$: parity.
For $k=2$, $H = \frac1{\sqrt2}\begin{psmallmatrix}1&1\\1&-1\end{psmallmatrix}$,
$T = \operatorname{diag}(1,i)$: $HT = \frac1{\sqrt2}\begin{psmallmatrix}1&i\\1&-i\end{psmallmatrix}$,
$(HT)^2 = \frac{1+i}{2}\begin{psmallmatrix}1&1\\-i&i\end{psmallmatrix}$,
$(HT)^3 = \frac{1+i}{\sqrt2}I = e^{i\pi/4}I$, while $S^2 = $ parity $= I$ at $k=2$.
So $(ST)^3 = e^{i\pi/4}S^2$: the relation holds projectively, as the fit's
group-law check reports ($9.5\times10^{-16}$). (With the $c/24$ normalization
of the affine $T$, P64, the relation becomes exact.)
\end{solution}

\begin{solution}{33}
$(n+j/k)^{\mathsf T}B(n+j/k) = n^{\mathsf T}Bn + \tfrac2k n^{\mathsf T}Bj + \tfrac1{k^2}j^{\mathsf T}Bj$.
Times $\pi ik$: $\pi ik\,n^{\mathsf T}Bn = \pi ik\sum_iB_{ii}n_i^2 + 2\pi ik\sum_{i<l}B_{il}n_in_l$,
every term in $2\pi i\Z$ for $k$ even; $2\pi i\,n^{\mathsf T}Bj\in2\pi i\Z$; the
remainder is $e^{\pi ij^{\mathsf T}Bj/k}$. For $B = \begin{psmallmatrix}0&1\\1&0\end{psmallmatrix}$,
$j^{\mathsf T}Bj = 2j_1j_2$, so the phase is $e^{2\pi ij_1j_2/k} = (-1)^{j_1j_2}$ at
$k=2$: $\operatorname{diag}(1,1,1,-1)$, controlled-$Z$.
\end{solution}

\begin{solution}{34}
For block-diagonal $\mathsf M$ the Siegel action, the map on $z$ and the automorphy
factor all split, and by P24 the basis splits, so the transformation law is
the tensor product of the two laws and $\rho(\mathsf M_1\oplus\mathsf M_2) = \rho(\mathsf M_1)\otimes\rho(\mathsf M_2)$.
$CZ = |0\rangle\langle0|\otimes I + |1\rangle\langle1|\otimes Z$ is a sum of two
products and not one: operator Schmidt rank $2$.
\end{solution}

\begin{solution}{35}
For fixed $\tau$ the functions $\theta_l(\cdot;\tau)$ are linearly independent
(distinct clock eigenvalues, P23), so the coefficients $\rho_{jl}$ in the law
are uniquely determined for each $\tau$. The derivations of P30--P31 produce
them as roots of unity over $\sqrt k$ with no $\tau$ anywhere, so the unique
coefficients are the same at every $\tau$. Equivalently $\rho_k$ is a
representation of a finite group built from the finite Heisenberg group,
which has no continuous parameter. Measured: $9\times10^{-16}$.
\end{solution}

\begin{solution}{36}
Let $\omega_1,\dots,\omega_g$ be the holomorphic forms with $\oint_{A_i}\omega_j = \delta_{ij}$,
$\oint_{B_i}\omega_j = \Omega_{ij}$. A period map acting on periods by
$a' = Pa+Qb$, $b' = Ra+Sb$ corresponds to the new cycles $A' = PA+QB$,
$B' = RA+SB$ (as chains), so $\oint_{A'_i}\omega_j = (P+Q\Omega)_{ij}$ and
$\oint_{B'_i}\omega_j = (R+S\Omega)_{ij}$. Renormalizing $\omega' = \omega(P+Q\Omega)^{-1}$
makes the $A'$-periods the identity and the $B'$-periods
$\Omega' = (R+S\Omega)(P+Q\Omega)^{-1}$. Matching with $(A\Omega+B)(C\Omega+D)^{-1}$
gives $A = S$, $B = R$, $C = Q$, $D = P$. For $z$: $z' = Ra+Sb-\Omega'(Pa+Qb)
= (S-\Omega'Q)b+(R-\Omega'P)a$, and from $\Omega'(P+Q\Omega) = R+S\Omega$,
$R-\Omega'P = -(S-\Omega'Q)\Omega$, so $z' = (S-\Omega'Q)(b-\Omega a)$. The
standard identity $(C\Omega+D)^{-\mathsf T} = A-\Omega'C$ for symplectic
$\mathsf M$ (equivalently $M_{\rm per}$ symplectic) reads $S-\Omega'Q = (P+Q\Omega)^{-\mathsf T}$.
Genus one: $\tau' = (R+S\tau)/(P+Q\tau) = (M_{21}+M_{22}\tau)/(M_{11}+M_{12}\tau)$,
which is the relation F3 measures between the flipped far torus's own modulus
in the relabelled marking and the near modulus, to $8\times10^{-16}$. Feeding
$M_{\rm per}$ itself into the Siegel formula would use $(P\tau+Q)/(R\tau+S)$,
which is wrong.
\end{solution}

\begin{solution}{37}
The $Z_r$ commute and are diagonal in the theta basis; $|0\rangle$ is a joint
eigenvector with all eigenvalues $1$. Its orbit $\{X^p|0\rangle = |p\rangle\}$ runs
over the basis, and $Z^q$ acts on $|p\rangle$ by the character $\omega^{q\cdot p}$,
so the $d$ vectors carry the $d$ distinct joint eigencharacters: any invariant
subspace containing one of them contains its whole shift orbit, hence
everything, so the representation is irreducible. A compatible scalar lift of
a symplectic relabelling sends the generators to new operators $X'^pZ'^q$
obeying the same relations with the same central phases, so it defines a second
irreducible representation on the same space. Both are irreducible
representations of the finite Heisenberg group with the same central
character, so (Stone--von Neumann for finite groups, or directly: build the
new joint eigenvector and map orbits to orbits) there is an intertwiner $U$
with $U D_{p,q}U^{-1} = D'_{p,q}$. Since the $D_{p,q}$ span
$\operatorname{End}(\mathcal Q_k)$, any two intertwiners differ by an operator
commuting with everything, a scalar. If the lift preserves adjoints, then
$U^\dagger U$ commutes with every $D_{p,q}$, so it is a positive scalar;
rescaling makes $U$ unitary, unique up to a phase. A composite of compatible
lifts has an intertwiner that is a composite of intertwiners, unique up to
phase, so the lifts compose projectively.
\end{solution}

\section*{Orientation reversal and coherent states (Sections~14--15)}

\begin{solution}{38}
$\overline{\theta_j(z;\Omega)} = \sum_n\exp(-\pi ik(n+j/k)^{\mathsf T}\bar\Omega(n+j/k) - 2\pi ik(n+j/k)^{\mathsf T}\bar z)
= \sum_n\exp(\pi ik(n+j/k)^{\mathsf T}(-\bar\Omega)(n+j/k) + 2\pi ik(n+j/k)^{\mathsf T}(-\bar z)) = \theta_j(-\bar z;-\bar\Omega)$.
$R_{\rm per}:(a,b)\mapsto(a,-b)$ sends $z = b-\Omega a$ to $-b-\Omega a$, which
is $-\bar z$ for the polarization $-\bar\Omega$; so $R_{\rm per}$ acts on the
value column by complex conjugation $K$. $\sigma^{\mathsf T}J\sigma = -J$
(P6), and $N = \sigma R_{\rm per} = \begin{psmallmatrix}0&-1\\1&0\end{psmallmatrix} = S$
satisfies $N^{\mathsf T}JN = J$. Its Siegel matrix, by P36 with
$P = 0,Q = -1,R = 1,S = 0$, is $\mathsf N = \begin{psmallmatrix}0&1\\-1&0\end{psmallmatrix} = S^{-1}$.
Hence the swap acts on values by $\rho_2(S^{-1})K = HK$ up to a phase
($H^{-1} = H$): Hadamard composed with conjugation, antiunitary.
\end{solution}

\begin{solution}{39}
Conjugating the law: $\overline{F_{\Omega'}(z')} = \bar f\,\bar\rho\,\overline{F_\Omega(z)}$.
Norms: $\|\overline{F_{\Omega'}(z')}\| = |f|\,\|\bar\rho\overline{F_\Omega(z)}\| = |f|\,\|F_\Omega(z)\|$
because $\bar\rho$ is unitary when $\rho$ is. Dividing,
$c_{\Omega'}(z') = (\bar f/|f|)\bar\rho\,c_\Omega(z)$. The coherent state is the
Riesz representative of evaluation, so it transforms by the conjugate
representation: applying $\rho$ instead of $\bar\rho$ gives a vector that is
not the coherent state at the transported point (for the Hadamard, $\bar\rho = \rho$,
which hides the error at level two for $S$, but $\bar\rho(T) = \operatorname{diag}(1,-i)\ne\rho(T)$).
Both $\rho$ and $\bar\rho$ are faithful unitary representations of the same
finite group, and their images span the same matrix algebra (the complex span
of the translations is $\operatorname{End}(\mathcal Q_k)$ in either), so no
completeness is lost by the choice; only consistency between fitted and
applied representations matters.
\end{solution}

\begin{solution}{40}
On a genus-one curve, Riemann--Roch gives $h^0(D)-h^0(K-D) = \deg D$ and $K = 0$,
so $h^0(D) = \deg D$ for $\deg D>0$. Thus $h^0(L^2) = 2$ and, for any point $p$,
$h^0(L^2(-p)) = 1$: the sections vanishing at $p$ form a line, so both basis
sections cannot vanish at $p$ simultaneously (else $h^0(L^2(-p)) = 2$). Hence
$[\theta_0:\theta_1]$ is a well-defined holomorphic map $E\to\mathbb{CP}^1$.
It is nonconstant (the two sections are independent), so by the open mapping
theorem its image is open, and by compactness closed: all of $\mathbb{CP}^1$.
Its degree is the number of zeros of a generic section $\alpha\theta_0+\beta\theta_1$,
which is $\deg L^2 = 2$ (P22); so every ray has two preimages generically.
Composing with the antiholomorphic bijection $[u:v]\mapsto[\bar u:\bar v]$ of
$\mathbb{CP}^1$ preserves surjectivity, so $z\mapsto c_\tau(z)$ is onto. For
a qutrit the image of $E$ in $\mathbb{CP}^2$ is a curve (real dimension $2$)
inside a manifold of real dimension $4$: a proper subset. In general the
coherent family has complex dimension at most $g$ against $k^g-1$.
\end{solution}

\begin{solution}{41}
By Proposition~11.1, $\theta_{(j_1,j_2)}(z;\Omega) = \theta_{j_1}(z_1;\tau_1)\theta_{j_2}(z_2;\tau_2)$,
so the value column is the Kronecker product $F_{\tau_1}(z_1)\otimes F_{\tau_2}(z_2)$;
conjugating and normalizing preserves the product structure, and a product
vector has one Schmidt coefficient, entropy $0$ (measured $5\times10^{-30}$).
With $\Omega_{12}\ne0$ the lattice sum does not factor (P24), so the
$2\times2$ coefficient matrix generically has rank two. The register's tensor
labels are a declared identification of $\mathcal Q_2(\Omega)$ with
$\mathcal Q_2\otimes\mathcal Q_2$ (through the identity monodromy); the
entanglement is a property of the encoded state at the holonomy point, which
depends on the polarization $\Omega$, not of the operator $\rho_2(I) = I$.
\end{solution}

\section*{Seams, factorization and the Clifford limit (Sections~16--17)}

\begin{solution}{42}
$U = W_1$, $V = W_2$, $U\cap V\simeq S^2$:
$H^0(W_1)\oplus H^0(W_2)\to H^0(S^2)$ is onto, and $H^1(S^2)=0$, so
$H^1(W)\to H^1(W_1)\oplus H^1(W_2)$ is an isomorphism. Every harmonic form is a
pair with one member on each collar, so the restriction to the four tori is
$M_{01}\oplus M_{23}$ and $\rho = \rho(M_{01})\otimes\rho(M_{23})$
(Proposition~13.3). The off-block norm tests that the isomorphism commutes
with restriction to the boundary: no harmonic form has periods on both
collars' tori.
\end{solution}

\begin{solution}{43}
$U\cap V\simeq S^1$, so $H^1(W)\to H^1(W_1)\oplus H^1(W_2)\to H^1(S^1)=\Z$ with
kernel $\{(\alpha_1,\alpha_2):\alpha_1(a_1) = \alpha_2(a_2)\}$; $H^0(U)\oplus H^0(V)\to H^0(S^1)$
is onto, so nothing new enters from degree zero and $H^1(W)$ is that kernel,
of rank $3$. On the boundary the restriction subspace obeys $p_1(a_1) = p_2(a_2)$:
the two tori's periods are coupled. The simplest symplectic matrix with a
coupled graph is the shear $\begin{psmallmatrix}I&B\\0&I\end{psmallmatrix}$,
$B = \begin{psmallmatrix}0&1\\1&0\end{psmallmatrix}$, whose $\rho_2$ is controlled-$Z$
(P33) of Schmidt rank $2$ (P34).
\end{solution}

\begin{solution}{44}
Cut the handle in its middle disk: $U\simeq W_1$, $V\simeq W_2$, $U\cap V\simeq D^2$,
contractible. $H^0(U)\oplus H^0(V)\to H^0(D^2)$ is onto and $H^1(D^2)=0$, so
$H^1(W) = H^1(W_1)\oplus H^1(W_2)$ exactly as for the sphere: a $1$-handle
between two components is a path, $W_1\cup\text{handle}\cup W_2\simeq W_1\vee W_2$.
An annulus has $H^1(S^1)=\Z$, which is what produces the coupling of P43.
The tube host is a $1$-handle, so its relation is block-diagonal: this is
topology, and G2 confirms it ($M = I_4$, off-block norm $2.6\times10^{-15}$).
The far boundary, however, is one connected genus-two surface, whose period
matrix is defined by its harmonic $1$-forms for the boundary metric; the
neck lets a harmonic form with periods on one torus leak into the other, so
$\Omega_{12}\ne0$ with a size set by the neck's waist and length. That is a
metric statement about the polarization $\Omega$ of one surface, not about
the cobordism's relation: as the neck pinches, $\Omega_{12}\to0$ and the
diagonal approaches the two tori's moduli (G5, the separating degeneration).
\end{solution}

\begin{solution}{45}
$|{+}{+}\rangle = \tfrac12(|00\rangle+|01\rangle+|10\rangle+|11\rangle)$ and $CZ$ flips
the sign of $|11\rangle$: $CZ|{+}{+}\rangle = \tfrac12(|00\rangle+|01\rangle+|10\rangle-|11\rangle)$.
Its coefficient matrix $\tfrac12\begin{psmallmatrix}1&1\\1&-1\end{psmallmatrix}$
is $\tfrac1{\sqrt2}H$, with two equal singular values $\tfrac1{\sqrt2}$: the
reduced state is $I/2$, one bit of entanglement. $\mathrm{SWAP}(|a\rangle\otimes|b\rangle) = |b\rangle\otimes|a\rangle$
is a product for every product input, yet $\mathrm{SWAP} = \tfrac12\sum_{P\in\{I,X,Y,Z\}}P\otimes P$
has operator Schmidt rank four. The witness: an operator is entangling if
some product state is mapped to a state with nonzero reduced entropy; E5
evaluates the geometric gate on all $36$ product stabilizer states and finds
$16$ of them mapped to Bell pairs (exactly one bit). Operator Schmidt rank
measures operator-vector correlations (the Choi state), which SWAP has
maximally while entangling nothing; Proposition~13.3 only shows that rank
one is \emph{not} entangling.
\end{solution}

\begin{solution}{46}
$S$ sends $\tau\mapsto-1/\tau$ and $z\mapsto z/\tau$, exchanging the roles of the
lattice directions $1$ and $\tau$; the translation by $\tau/k$ ($\mathsf X$)
becomes the translation by $1/k$ ($\mathsf Z$) up to the automorphy phases, so
$\rho(S)\mathsf X\rho(S)^{-1}\propto\mathsf Z$. Translations by $k$-torsion points
form the Heisenberg--Weyl group, and a unitary conjugating it to itself lies in
its normalizer, the Clifford group; for $k=2$ the group generated by $H$ and
$\operatorname{diag}(1,i)$ is the whole single-qubit Clifford group (order $24$
modulo phases). Gottesman--Knill: a Clifford circuit acting on a stabilizer
state can be tracked by updating $2n$ generators of its stabilizer, polynomial
in $n$, so the theta register alone yields nothing beyond classical simulation.
This does not contradict Theorem~12.2: the complex \emph{span} of the
translations is the full matrix algebra, while the \emph{group} generated by
the modular gates is finite.
\end{solution}

\begin{solution}{47}
$\Sp(2,\mathbb F_2) = \SL(2,\mathbb F_2)\cong S_3$, so the local Cliffords reduce
into $S_3\times S_3$, and with the swap into $S_3\wr\Z/2$, of order
$6\cdot6\cdot2 = 72$. An element of order five in $S_6$ is a $5$-cycle; $5\nmid72$,
so no conjugate of $S_3\wr\Z/2$ contains it. A local Clifford $A\otimes B$ (or
$\mathrm{SWAP}\cdot A\otimes B$) has symplectic image in that subgroup, so an
element of order five modulo $2$ cannot be one: its entangling capacity is
forced by its order. For $M$: expanding along the third row,
$\det M = 1$; $J = \begin{psmallmatrix}0&I\\-I&0\end{psmallmatrix}$ and a direct
multiplication gives $M^{\mathsf T}JM = J$; $M^2 = \begin{psmallmatrix}0&-1&1&1\\0&-1&0&1\\-1&0&0&0\\1&-1&0&0\end{psmallmatrix}$
and continuing, $M^5 = I$ (the characteristic polynomial is
$x^4+x^3+x^2+x+1$, the fifth cyclotomic polynomial). Modulo $2$, $M\not\equiv I$
and $M^5\equiv I$, so the order is five. The measured gate has operator
Schmidt rank four and sends $16$ of $36$ product stabilizer states to Bell
pairs (E4, E5), as the argument predicts.
\end{solution}

\section*{From a Lie algebra to its root system; the Weyl group (Sections~18--19)}

\begin{solution}{48}
Let $v$ be a highest-weight vector: $ev = 0$, $hv = \lambda v$. From
$[h,f] = -2f$, $h(f^nv) = (\lambda-2n)f^nv$. From $[e,f] = h$ and
induction, $ef^nv = n(\lambda-n+1)f^{n-1}v$: indeed
$ef^nv = fef^{n-1}v + hf^{n-1}v = [(n-1)(\lambda-n+2)+(\lambda-2n+2)]f^{n-1}v
= n(\lambda-n+1)f^{n-1}v$. In finite dimension the string terminates: let
$N$ be least with $f^{N+1}v = 0$, $f^Nv\ne0$; then
$0 = ef^{N+1}v = (N+1)(\lambda-N)f^Nv$ forces $\lambda = N$. The span of
$v,fv,\dots,f^\lambda v$ is invariant under $e,f,h$, so by irreducibility it
is the module: eigenvalues $\lambda,\lambda-2,\dots,-\lambda$, each once,
dimension $\lambda+1$. Identification: $[h,e] = 2e$ says $\alpha(h) = 2$;
the coroot is the element $\alpha^\vee = h$ of $\mathfrak h$ with
$\langle\mu,\alpha^\vee\rangle = \mu(h)$; $\omega(h) = 1$ so
$\alpha = 2\omega$; $P = \Z\omega$ (the $h$-eigenvalues are integers),
$Q = \Z\alpha = 2\Z\omega$, $P/Q = \Z/2$; $s_\alpha(\omega) = \omega-\langle\omega,\alpha^\vee\rangle\alpha = \omega-2\omega = -\omega$,
so $W = \{\pm1\}$. $|\omega|^2 = |\alpha|^2/4 = \tfrac12$.
\end{solution}

\begin{solution}{49}
$[\operatorname{diag}(a),E_{ij}] = (a_i-a_j)E_{ij}$, so the roots are
$\varepsilon_i-\varepsilon_j$ with root vectors $E_{ij}$. With
$\langle\varepsilon_i,\varepsilon_j\rangle = \delta_{ij}-\tfrac13$:
$|\varepsilon_i-\varepsilon_j|^2 = 2$ and
$\langle\alpha_1,\alpha_2\rangle = \langle\varepsilon_1,\varepsilon_2\rangle-\langle\varepsilon_1,\varepsilon_3\rangle-\langle\varepsilon_2,\varepsilon_2\rangle+\langle\varepsilon_2,\varepsilon_3\rangle
= -\tfrac13+\tfrac13-\tfrac23-\tfrac13 = -1$. Cartan matrix
$A_{ij} = 2\langle\alpha_i,\alpha_j\rangle/|\alpha_j|^2 = \begin{psmallmatrix}2&-1\\-1&2\end{psmallmatrix}$;
$\cos\angle(\alpha_1,\alpha_2) = -1/2$, so $120^\circ$. The positive roots
are $\alpha_1,\alpha_2,\alpha_1+\alpha_2 = \varepsilon_1-\varepsilon_3 = \theta$.
Fundamental weights: $\omega_1 = a\alpha_1+b\alpha_2$ with
$\langle\omega_1,\alpha_1\rangle = 2a-b = 1$, $\langle\omega_1,\alpha_2\rangle = -a+2b = 0$
gives $\omega_1 = \tfrac13(2\alpha_1+\alpha_2)$, and by symmetry
$\omega_2 = \tfrac13(\alpha_1+2\alpha_2)$; $P/Q\cong\Z/3$ generated by
$\omega_1$ ($3\omega_1\in Q$, $2\omega_1 = \omega_2+\alpha_1$).
$\rho = \omega_1+\omega_2 = \alpha_1+\alpha_2 = \theta$;
$h^\vee = \langle\rho,\theta\rangle+1 = 2+1 = 3$. The reflection
$s_{\varepsilon_i-\varepsilon_j}$ exchanges $\varepsilon_i$ and
$\varepsilon_j$ (e.g.\ $s_{\alpha_1}\varepsilon_1 = \varepsilon_1-(\tfrac23+\tfrac13)(\varepsilon_1-\varepsilon_2) = \varepsilon_2$),
so $W = S_3$ permuting the $\varepsilon_i$, the symmetry group of an
equilateral triangle, $D_3$. $P^+$ is the $60^\circ$ cone spanned by
$\omega_1,\omega_2$ ($|\omega_i|^2 = \tfrac23$, $\langle\omega_1,\omega_2\rangle = \tfrac13$).
\end{solution}

\begin{solution}{50}
$A_1\times A_1$: $A = \operatorname{Gram} = \operatorname{diag}(2,2)$.
$A_2$: both $\begin{psmallmatrix}2&-1\\-1&2\end{psmallmatrix}$. $B_2$
($\alpha_1$ long, $\alpha_2$ short, $|\alpha_2|^2 = 1$, angle $135^\circ$):
$\langle\alpha_1,\alpha_2\rangle = -1$, Gram
$\begin{psmallmatrix}2&-1\\-1&1\end{psmallmatrix}$,
$A = \begin{psmallmatrix}2&-2\\-1&2\end{psmallmatrix}$. $G_2$ ($\alpha_2$
short, $|\alpha_2|^2 = \tfrac23$, angle $150^\circ$):
$\langle\alpha_1,\alpha_2\rangle = -1$, Gram
$\begin{psmallmatrix}2&-1\\-1&2/3\end{psmallmatrix}$,
$A = \begin{psmallmatrix}2&-3\\-1&2\end{psmallmatrix}$. Lattices: in $B_2$
the short roots $\alpha_2$ and $\alpha_1+\alpha_2$ are orthogonal of
length $1$ ($\langle\alpha_2,\alpha_1+\alpha_2\rangle = -1+1 = 0$,
$|\alpha_1+\alpha_2|^2 = 2-2+1 = 1$) and generate $Q$, since
$\alpha_1 = (\alpha_1+\alpha_2)-\alpha_2$: the unit square lattice; the
$A_1\times A_1$ root lattice is the square lattice of side $\sqrt2$. In $G_2$
the short roots $\pm\alpha_2,\pm(\alpha_1+\alpha_2),\pm(\alpha_1+2\alpha_2)$
have length $\sqrt{2/3}$ and $\langle\alpha_2,\alpha_1+\alpha_2\rangle = -1+\tfrac23 = -\tfrac13$,
i.e.\ $120^\circ$: they form an $A_2$ system and generate $Q$
($\alpha_1 = (\alpha_1+\alpha_2)-\alpha_2$), a hexagonal lattice, the $A_2$
lattice scaled by $1/\sqrt3$. Weyl groups: reflections in the two long
roots of $B_2$ (which are orthogonal) generate $D_2 = W(A_1\times A_1)$,
index two in $D_4$; the reflection in a short root is the missing
generator. The long roots of $G_2$, $\pm\alpha_1,\pm(\alpha_1+3\alpha_2),\pm(2\alpha_1+3\alpha_2)$,
form an $A_2$ system whose reflections generate $D_3 = W(A_2)$, index two
in $D_6$; again a short-root reflection is missing.
\end{solution}

\begin{solution}{51}
With $\alpha_1,\alpha_2$ the sides from the vertex $0$ and $\alpha_1-\alpha_2$
the third side, $\alpha_1\cdot\alpha_2 = (2+2-6)/2 = -1$, giving $A$. The
squared norm of $m\alpha_1+n\alpha_2$ is $2m^2-2mn+2n^2 = 2(m^2-mn+n^2)$;
$m^2-mn+n^2 = \tfrac34n^2+(m-\tfrac n2)^2\ge1$ for $(m,n)\ne0$ integers, with
equality exactly at $(\pm1,0),(0,\pm1),\pm(1,1)$: six vectors of norm $2$ at
mutual angles multiples of $60^\circ$, the root system $A_2$. Since
$|\alpha_i|^2 = 2$, $s_{\alpha_i}(v) = v-(v\cdot\alpha_i)\alpha_i$, and
$v\cdot\alpha_i\in\Z$ for lattice $v$ (integral Gram), so each reflection maps
$Q$ to itself; $s_1s_2$ is the rotation by $120^\circ$, of order three, and
$\langle s_1,s_2\rangle$ is the dihedral group of order six, $S_3$. The dual
lattice $\{v: v\cdot Q\subseteq\Z\}$ in root coordinates is $A^{-1}\Z^2$, and
$[P:Q] = \det A = 3$. Everything used only the three lengths (the Gram matrix
is determined by them) and integrality; no Lie algebra was input. For a
boundary triangulation's harmonic Gram to equal $A$ the triangulation would
need a harmonic basis whose mass pairing is $A$, i.e.\ a $D_3$-symmetric
metric and marking; a generic triangulation has neither, which is why the
reading places the root lattice on an internal fiber (Section~32) rather
than on the boundary.
\end{solution}

\begin{solution}{52}
$\mathfrak{su}(2)$: $W = \{1,s\}$, $\varepsilon(s) = -1$,
$s(\lambda+\rho) = -(\lambda+1)\omega$, so with $x = e^{\omega}$ the
numerator is $x^{\lambda+1}-x^{-\lambda-1}$ and the denominator
$x-x^{-1}$; the quotient is
$x^{\lambda}(1-x^{-2(\lambda+1)})/(1-x^{-2}) = x^{\lambda}\sum_{i=0}^{\lambda}x^{-2i}$.
$\mathfrak{su}(3)$: $\lambda = a\omega_1+b\omega_2$, $\rho = \omega_1+\omega_2$,
positive coroots $\alpha_1,\alpha_2,\theta$;
$\langle\lambda+\rho,\alpha_1\rangle = a+1$,
$\langle\lambda+\rho,\alpha_2\rangle = b+1$,
$\langle\lambda+\rho,\theta\rangle = a+b+2$, against $1,1,2$ for $\rho$:
$\dim = (a+1)(b+1)(a+b+2)/2$. Checks: $2\cdot1\cdot3/2 = 3$,
$2\cdot2\cdot4/2 = 8$, $4\cdot1\cdot5/2 = 10$, $3\cdot3\cdot6/2 = 27$.
\end{solution}

\begin{solution}{53}
Proof: $\big(\sum_w\varepsilon(w)e^{w(\lambda+\rho)}\big)\big(\sum_\nu m_\mu(\nu)e^{\nu}\big)
= \sum_w\varepsilon(w)\sum_{\nu}m_\mu(\nu)e^{w(\lambda+\rho)+\nu}$;
substitute $\nu = w\nu'$ (a bijection of $P$ with $m_\mu(w\nu') = m_\mu(\nu')$)
to get $\sum_{\nu'}m_\mu(\nu')\sum_w\varepsilon(w)e^{w(\lambda+\rho+\nu')}$.
If $\lambda+\rho+\nu'$ is fixed by a reflection $s$, pairing $w$ with $ws$
cancels the inner sum; otherwise $\lambda+\rho+\nu' = w_{\nu'}^{-1}\kappa$ with
$\kappa$ strictly dominant, and the inner sum is
$\varepsilon(w_{\nu'})\sum_w\varepsilon(w)e^{w\kappa}$, the numerator of
$\operatorname{ch}V_{\kappa-\rho}$. Divide by the Weyl denominator.
$\mathfrak{su}(2)$, $V_1\otimes V_\lambda$: $\nu = \pm1$ gives
$V_{\lambda+1}$ and, for $\lambda\ge1$, $V_{\lambda-1}$; at $\lambda = 0$,
$\lambda+\nu+\rho = 0$ is on the wall and is killed: $V_1\otimes V_0 = V_1$.
$V_2\otimes V_0$: weights of $V_2$ are $2,0,-2$; $\nu = 2$ gives $V_2$,
$\nu = 0$ gives $V_0$, $\nu = -2$ gives $\lambda+\nu+\rho = -1$, folded by
$s$ to $+1 = \rho+0$ with sign $-1$: $-V_0$. Total $V_2+V_0-V_0 = V_2$.
$\mathfrak{su}(3)$, $\mathbf3\otimes\mathbf3$ ($\lambda = (1,0)$, weights
of $\mathbf3$: $(1,0),(-1,1),(0,-1)$ in the $\omega$ basis):
$(2,0) = \mathbf6$; $(0,1) = \bar{\mathbf3}$; $(1,-1)+\rho = (2,0)$ lies on
the wall $\langle\cdot,\alpha_2^\vee\rangle = 0$ and is killed:
$\mathbf3\otimes\mathbf3 = \mathbf6\oplus\bar{\mathbf3}$ ($9 = 6+3$).
$\mathbf3\otimes\bar{\mathbf3}$ ($\lambda = (0,1)$): $(1,1) = \mathbf8$;
$(-1,2)+\rho = (0,3)$ on a wall, killed; $(0,0) = \mathbf1$:
$\mathbf8\oplus\mathbf1$.
\end{solution}

\begin{solution}{54}
$A_1$: $e^{\omega}-e^{-\omega} = e^{\omega}(1-e^{-2\omega}) = e^{\rho}(1-e^{-\alpha})$.
$A_2$: $e^{\rho}(1-e^{-\alpha_1})(1-e^{-\alpha_2})(1-e^{-\alpha_1-\alpha_2})$
expands to eight terms, $e^{\rho}[1-e^{-\alpha_1}-e^{-\alpha_2}-e^{-\alpha_1-\alpha_2}
+e^{-\alpha_1-\alpha_2}+e^{-2\alpha_1-\alpha_2}+e^{-\alpha_1-2\alpha_2}-e^{-2\alpha_1-2\alpha_2}]$;
the two $e^{\rho-\alpha_1-\alpha_2}$ terms cancel, leaving six. They are
the $W$-orbit of $\rho$ with signs: $\rho$ ($1$, $+$);
$s_1\rho = \rho-\alpha_1$ ($-$); $s_2\rho = \rho-\alpha_2$ ($-$);
$s_1s_2\rho = \rho-\alpha_1-s_1\alpha_2 = \rho-2\alpha_1-\alpha_2$ ($+$);
$s_2s_1\rho = \rho-\alpha_1-2\alpha_2$ ($+$);
$s_1s_2s_1\rho = -\rho = \rho-2\alpha_1-2\alpha_2$ ($-$).
\end{solution}

\section*{Affine algebras and theta functions (Section~20)}

\begin{solution}{55}
$wt_\gamma w^{-1}(\mu) = w(w^{-1}\mu+K\gamma) = \mu+Kw\gamma = t_{w\gamma}(\mu)$,
so $t(Q^\vee)$ is normalized by $W$ ($W$ preserves $Q^\vee$); a translation
fixes no point unless $\gamma = 0$ while $W$ fixes $0$, so
$W\cap t(Q^\vee) = 1$ and the group generated is $W\ltimes t(Q^\vee)$.
Since $\theta$ is long, $\theta^\vee = \theta$ and
$s_0(\mu) = s_\theta(\mu)+K\theta = \mu-(\langle\mu,\theta\rangle-K)\theta$,
the reflection in the affine hyperplane $\langle\mu,\theta\rangle = K$
along $\theta$: it fixes that hyperplane pointwise and squares to the
identity. $\mathfrak{su}(2)$, $K = k+2$, $\theta = \alpha = 2\omega$:
$s_0(m\omega) = -m\omega+2K\omega = (2(k+2)-m)\omega$, fixing $m = k+2$.
Alcove: the dominant chamber is a simplicial cone with $r$ walls, and
$\theta^\vee = \sum_ia_i^\vee\alpha_i^\vee$ with all $a_i^\vee>0$, so the
half-space $\langle\mu,\theta^\vee\rangle\le K$ cuts it to a bounded
simplex. Interior lattice points satisfy $\langle\mu,\alpha_i^\vee\rangle\ge1$
and $\langle\mu,\theta^\vee\rangle\le K-1$; with $\mu = \lambda+\rho$ these
read $\langle\lambda,\alpha_i^\vee\rangle\ge0$ and
$\langle\lambda,\theta^\vee\rangle\le K-1-\langle\rho,\theta^\vee\rangle = K-h^\vee = k$:
exactly $P_k^+ + \rho$.
\end{solution}

\begin{solution}{56}
$e^{t_\gamma(\hat\mu)} = e^{\mu+K\gamma}\,e^{-(\langle\mu,\gamma\rangle+\frac K2|\gamma|^2)\delta}
= e^{\mu+K\gamma}\,q^{\langle\mu,\gamma\rangle+\frac K2|\gamma|^2}$, and
$\tfrac K2|\tfrac\mu K+\gamma|^2 = \tfrac{|\mu|^2}{2K}+\langle\mu,\gamma\rangle+\tfrac K2|\gamma|^2$,
so $q^{\langle\mu,\gamma\rangle+\frac K2|\gamma|^2} = q^{-|\mu|^2/2K}q^{\frac K2|\frac\mu K+\gamma|^2}$;
with $e^{\mu+K\gamma} = e^{K(\frac\mu K+\gamma)}\mapsto e^{2\pi iK\langle\frac\mu K+\gamma,z\rangle}$
the sum over $\gamma$ is $q^{-|\mu|^2/2K}\Theta^{(K)}_\mu$. Every
$\hat w\in\hat W$ is uniquely $t_\gamma w$; $\varepsilon(t_\gamma) = +1$
because $t_{\theta^\vee} = s_0s_\theta$ is a product of two reflections and
the $t_{w\theta^\vee}$ generate $t(Q^\vee)$. Hence
$\sum_{\hat w}\varepsilon(\hat w)e^{\hat w(\hat\lambda+\hat\rho)}
= \sum_{w\in W}\varepsilon(w)\sum_\gamma e^{t_\gamma(w(\hat\lambda+\hat\rho))}
= q^{-|\lambda+\rho|^2/2K}\sum_w\varepsilon(w)\Theta^{(K)}_{w(\lambda+\rho)}$,
the prefactor being common because $W$ preserves norms. The lattice enters
through $Q^\vee$ and the quadratic form $|\cdot|^2$ (the theta function);
the non-abelian structure enters only through $W$ and its sign, and
through $\rho$ and $h^\vee$ in the shift $\lambda\mapsto\lambda+\rho$,
$k\mapsto k+h^\vee$.
\end{solution}

\begin{solution}{57}
$h_\lambda-\tfrac c{24} = \dfrac{|\lambda+\rho|^2-|\rho|^2}{2(k+h^\vee)}-\dfrac{k\dim\mathfrak g}{24(k+h^\vee)}$.
The strange formula gives $|\rho|^2 = h^\vee\dim\mathfrak g/12$, so
$\dfrac{|\rho|^2}{2(k+h^\vee)}+\dfrac{k\dim\mathfrak g}{24(k+h^\vee)}
= \dfrac{(h^\vee+k)\dim\mathfrak g}{24(k+h^\vee)} = \dfrac{\dim\mathfrak g}{24}$.
$\mathfrak{su}(2)$: $\dim\mathfrak g = 3$, $|\rho|^2 = |\omega|^2 = \tfrac12$,
$h^\vee = 2$, and $\tfrac12/4 = \tfrac18 = \tfrac3{24}$ checks the strange
formula; $h_\lambda = \big(\tfrac{(\lambda+1)^2}{2}-\tfrac12\big)/2(k+2) = \lambda(\lambda+2)/4(k+2)$,
$T_\lambda = \exp2\pi i[(\lambda+1)^2/4(k+2)-\tfrac18]$. Central charges:
$\mathfrak{su}(2)_1$: $3/3 = 1$; $\mathfrak{su}(2)_2$: $6/4 = \tfrac32$;
$\mathfrak{su}(3)_1$: $8/4 = 2$. The $\mathfrak{su}(2)_1$ doublet:
$h_0 = 0$, $h_1 = 3/12 = \tfrac14$, so
$T = e^{-2\pi i/24}\operatorname{diag}(1,e^{2\pi i/4}) = e^{-\pi i/12}\operatorname{diag}(1,i)$,
which is $\rho_2(T) = \operatorname{diag}(e^{\pi ij^2/2}) = \operatorname{diag}(1,i)$
up to the overall phase.
\end{solution}

\begin{solution}{58}
$Q = \sqrt2\,\Z$, $\omega = 1/\sqrt2$, so $Q+j\omega = \{\sqrt2(n+\tfrac j2)\}$
and $\Theta_{Q+j\omega} = \sum_nq^{(n+j/2)^2}e^{2\pi i\sqrt2(n+j/2)z}$, which
is $\theta^{(2)}_j(z/\sqrt2;\tau) = \sum_n\exp(2\pi i(n+\tfrac j2)^2\tau+4\pi i(n+\tfrac j2)z/\sqrt2)$.
So $\chi_j = \theta^{(2)}_j/\eta$. Under $\tau\mapsto\tau+1$:
$\theta^{(2)}_j\mapsto e^{\pi ij^2/2}\theta^{(2)}_j$ and $\eta\mapsto e^{\pi i/12}\eta$,
so $T = e^{-\pi i/12}\operatorname{diag}(1,i)$, agreeing with P57. Under
$\tau\mapsto-1/\tau$: by P31, $\theta_j(z/\tau;-1/\tau) = (-i\tau)^{1/2}e^{\pi ikz^2/\tau}\cdot k^{-1/2}\sum_le^{-2\pi ijl/k}\theta_l(z;\tau)$
with $k = 2$, and $\eta(-1/\tau) = (-i\tau)^{1/2}\eta(\tau)$; the
$(-i\tau)^{1/2}$ cancels, leaving $S = \tfrac1{\sqrt2}\begin{psmallmatrix}1&1\\1&-1\end{psmallmatrix}$.
These are the Kac--Peterson matrices of P64 at $k = 1$. Fusion: Verlinde
with $S = H$ gives $N_{11}^{\ 0} = \sum_\sigma S_{1\sigma}^2 = 1$ and
$N_{11}^{\ 1} = \sum_\sigma S_{1\sigma}^3/S_{0\sigma} = 0$: $1\times1 = 0$,
addition in $\Z/2$, which is the composition
$\mathsf X^a\mathsf X^b = \mathsf X^{a+b}$ of characteristics in the
Heisenberg--Weyl group of P23.

$\mathfrak{su}(3)_1$: the cosets are $Q$, $Q+\omega_1$, $Q+\omega_2$ with
$\omega_2\equiv2\omega_1$; Poisson summation over the even lattice $Q$ (dual
$P$) gives $\Theta_{Q+a}(-1/\tau,z/\tau) = (-i\tau)e^{(\cdots)}|P/Q|^{-1/2}\sum_be^{-2\pi i\langle a,b\rangle}\Theta_{Q+b}(\tau,z)$,
and $\eta^2$ contributes $(-i\tau)$, so $S_{ab} = 3^{-1/2}e^{-2\pi i\langle\omega_a,\omega_b\rangle}$.
With $\langle\omega_1,\omega_1\rangle = \langle\omega_2,\omega_2\rangle = \tfrac23$,
$\langle\omega_1,\omega_2\rangle = \tfrac13$: $S_{11} = S_{22} = e^{2\pi i/3}/\sqrt3$,
$S_{12} = e^{4\pi i/3}/\sqrt3$, i.e.\ $S_{ab} = 3^{-1/2}e^{2\pi iab/3}$ for
$a,b\in\Z/3$, the Fourier transform of $\Z/3$. $T$: $h_a = |\omega_a|^2/2 = \tfrac13$
for $a = 1,2$ and $c = 2$, so $T = e^{-2\pi i/12}\operatorname{diag}(1,e^{2\pi i/3},e^{2\pi i/3})$.
Parity: on $Q(A_2)$, $|n\alpha_1+m\alpha_2|^2 = 2n^2-2nm+2m^2$ is even, so
$q^{|\gamma|^2/2}$ is invariant under $\tau\mapsto\tau+1$ term by term and
each coset picks up the constant $e^{\pi i|\omega_a|^2}$ (since
$|\omega_a+\gamma|^2 = |\omega_a|^2+2\langle\omega_a,\gamma\rangle+|\gamma|^2$
with $\langle\omega_a,\gamma\rangle\in\Z$). The rank-one register at level
$3$ is the lattice $\sqrt3\,\Z$ with $|\gamma|^2 = 3n^2$, odd for odd $n$:
$\tau\mapsto\tau+1$ multiplies the terms by $(-1)^{n}$, which depends on
$n$, and $\theta^{(3)}_j$ is not an eigenfunction --- the obstruction of
P30. Evenness of the lattice, not the rank or the level, is what makes $T$
diagonal.
\end{solution}

\section*{Rank two, the alcove, fusion and density (Sections~21--25)}

\begin{solution}{59}
Averaging any positive form over a finite $G\subset\GL(2,\Z)$ gives a
$G$-invariant one, so $G$ is a finite subgroup of $O(2)$: cyclic $C_n$ or
dihedral $D_n$. An element of $\GL(2,\Z)$ has integer trace, and a rotation by
$2\pi/n$ has trace $2\cos(2\pi/n)\in\Z$ only for $n\in\{1,2,3,4,6\}$. Up to
conjugacy: $C_1,C_2,C_3,C_4,C_6$ and $D_1,D_2,D_3,D_4,D_6$ (some with two
classes). Reflections: $D_2$ is generated by reflections in two orthogonal
lines --- the Weyl group of $A_1\times A_1$; $D_3$ by reflections at $60^\circ$
in the hexagonal lattice --- $W(A_2)$; $D_4$ by reflections in the axes and
diagonals of the square lattice --- $W(B_2)$; $D_6$ by the six reflections of
the hexagonal lattice --- $W(G_2)$.
\end{solution}

\begin{solution}{60}
Grid with the $(i,j)\!-\!(i{+}1,j{+}1)$ diagonal: the edge set is preserved by
$I$, the swap $(i,j)\mapsto(j,i)$, negation $(i,j)\mapsto(-i,-j)$, and their
product; the $90^\circ$ rotation sends the diagonal to the antidiagonal, which
is not an edge. Group $\{\pm I,\pm\sigma\}\cong D_2$, abelian. The triangular
lattice torus, with edges in three directions at $60^\circ$, is preserved by
every element of $D_6$: non-abelian. An automorphism of a finite complex
permutes finitely many simplices, so it has finite order, whereas
$T_B = \begin{psmallmatrix}1&1\\0&1\end{psmallmatrix}$ has infinite order in
$\GL(2,\Z)$. A flip pass replaces the far triangulation by a different one
(the grid in a sheared basis) and stacks one tetrahedron per flipped edge, so
the near and far faces are different triangulations of the same torus; the
far marking is read through the lattice map, giving $M = L^{\mathsf T}$, and
the quarter turn is the diagonal pass followed by the relabelling. This is a
change of complex, not a symmetry of one complex.
\end{solution}

\begin{solution}{61}
$\mathfrak{su}(2)$: $\theta^\vee$ is the coroot, $\langle\lambda\omega,\alpha^\vee\rangle = \lambda$,
so $\lambda\in\{0,\dots,k\}$: $k+1$ states. $\mathfrak{su}(3)$:
$\langle\lambda,\theta^\vee\rangle = \lambda_1+\lambda_2\le k$: $\binom{k+2}{2} = (k+1)(k+2)/2$
states; three at $k = 1$ and six at $k = 2$.
\end{solution}

\begin{solution}{62}
Units: fundamental weight $\omega$ with $|\omega|^2 = \tfrac12$, root $\alpha = 2\omega$,
$|\alpha|^2 = 2$, coroot $\alpha^\vee = \alpha$, $Q^\vee = \Z\alpha$, $\rho = \omega$,
$h^\vee = 2$. For $\mu = m\omega$ and $\gamma = n\alpha = 2n\omega$,
$\mu/K+\gamma = (m/K+2n)\omega$ and $|\mu/K+\gamma|^2 = \tfrac12(2n+m/K)^2 = 2(n+\tfrac{m}{2K})^2$,
so $\Theta^{(K)}_m = \sum_n\exp(2\pi iK(n+\tfrac m{2K})^2\tau + 2\pi iK(2n+\tfrac mK)\langle\omega,z\rangle)$,
which is the level-$2K$ theta function of characteristic $m\bmod 2K$ in the
variable $w = \langle\omega,z\rangle = \tfrac12\langle\alpha,z\rangle$. $W = \{\pm1\}$ acts by $m\mapsto-m$
with sign $-1$, so $A^{(K)}_m = \theta^{(2K)}_m - \theta^{(2K)}_{-m}$. The numerator
has $\mu = \lambda+\rho = (\lambda+1)\omega$ at level $K = k+h^\vee = k+2$, giving
$\theta^{(2(k+2))}_{\lambda+1}-\theta^{(2(k+2))}_{-(\lambda+1)}$; the denominator has
$\mu = \rho = \omega$ at level $h^\vee = 2$, giving $\theta^{(4)}_1-\theta^{(4)}_{-1}$.
Numerically, at $\tau = 0.13+i$ and $w = 0.17+0.06i$, both sides agree to
$7\times10^{-16}$ for $\lambda = 0,1$; and for $k = 1,2,3$ the quotient divided
by $q^{h_\lambda-c/24}$ tends to the finite character $\sum_{m\equiv\lambda(2),|m|\le\lambda}e^{2\pi imw}$
as $\operatorname{Im}\tau\to\infty$. At $w = 0$ the numerator and the
denominator both vanish, because $\theta_{-m}(0) = \theta_m(0)$ by evenness
($w = 0$ lies on the Weyl wall); the value there is a removable limit, which
is why the reading requires evaluation away from the walls or a controlled
limit.
\end{solution}

\begin{solution}{63}
With no positive roots the Weyl denominator is the empty product, $1$;
$W$ trivial makes $A_\mu = \Theta_\mu$; $\rho = 0$ and $h^\vee = 0$ leave
$\chi_\lambda = \Theta^{(k)}_\lambda$, the level-$k$ theta function of the charge
$\lambda\in\Z/k$: Definition~10.1 of the reading. In the non-abelian ratio the
denominator is $A^{(h^\vee)}_\rho$, the alternating sum at level $h^\vee$;
substituting $h^\vee = 0$ there gives a theta function at level zero, which is
not defined (the Gaussian factor disappears and the lattice sum diverges).
The abelian case is a separate construction, not a limit of the ratio.
\end{solution}

\begin{solution}{64}
Symmetry is visible. Unitarity: with $N = k+2$ and $m = \mu+1\in\{1,\dots,N-1\}$,
$\frac2N\sum_{m=1}^{N-1}\sin\frac{\pi am}N\sin\frac{\pi bm}N = \delta_{ab}$ for
$1\le a,b\le N-1$ (discrete sine transform), so $SS^{\mathsf T} = I$; $S$ is real
symmetric unitary, hence $S^2 = I$, the (trivial) charge conjugation of
$\mathfrak{su}(2)$. $k=1$: $S = \sqrt{\tfrac23}\cdot\tfrac{\sqrt3}2\begin{psmallmatrix}1&1\\1&-1\end{psmallmatrix}
= \tfrac1{\sqrt2}\begin{psmallmatrix}1&1\\1&-1\end{psmallmatrix}$, Hadamard;
$T = \operatorname{diag}(e^{-\pi i/12},e^{5\pi i/12}) = e^{-\pi i/12}\operatorname{diag}(1,i)$.
$k=2$: $S = \begin{psmallmatrix}\frac12&\frac1{\sqrt2}&\frac12\\ \frac1{\sqrt2}&0&-\frac1{\sqrt2}\\ \frac12&-\frac1{\sqrt2}&\frac12\end{psmallmatrix}$
(the Ising $S$-matrix), $T = \operatorname{diag}(e^{-\pi i/8},e^{\pi i/4},e^{7\pi i/8})$.
Exactness of $(ST)^3 = S^2$: at $k = 1$, $ST = e^{-\pi i/12}H\operatorname{diag}(1,i)$
and P32 gives $(H\operatorname{diag}(1,i))^3 = e^{\pi i/4}I$, so
$(ST)^3 = e^{-3\pi i/12}e^{\pi i/4}I = I = S^2$; the $c/24$ phase
$e^{-2\pi i c/24}$ with $c = 1$ exactly cancels the metaplectic phase. The
same holds at every $k$ (numerically to $6\times10^{-16}$ for $k\le4$): the
modular data of a rational theory is a true representation of $\SL(2,\Z)$
once $T$ carries $e^{-2\pi ic/24}$.
\end{solution}

\begin{solution}{65}
$k=2$, labels $0,1,2$. $N_{11}^{\ 0} = \sum_\sigma S_{1\sigma}^2 = 1$;
$N_{11}^{\ 1} = \sum_\sigma S_{1\sigma}^2S_{1\sigma}/S_{0\sigma}
= \tfrac12\cdot\tfrac{1}{\sqrt2}/\tfrac12 + 0 + \tfrac12\cdot(-\tfrac1{\sqrt2})/\tfrac12 = 0$;
$N_{11}^{\ 2} = \tfrac12\cdot\tfrac12/\tfrac12 + 0 + \tfrac12\cdot\tfrac12/\tfrac12 = 1$.
So $1\times1 = 0+2$ (Ising: $\sigma\times\sigma = 1+\psi$), matching the rule
$|\lambda-\mu|\le\nu\le\min(\lambda+\mu,2k-\lambda-\mu)$, $\lambda+\mu+\nu$ even.
Kac--Walton: the weights of $V_1$ are $\pm1$; $1+1 = 2$ and $1-1 = 0$ both lie
in the alcove, so $N_{11}^{\ 0} = N_{11}^{\ 2} = 1$ with no reflection. For
$2\times1$: $2+1 = 3$ lies on the affine wall $\lambda+1 = k+2$ and is fixed by
the affine reflection, contributing zero, while $2-1 = 1$ gives $N_{21}^{\ 1} = 1$
($\psi\times\sigma = \sigma$). Against Proposition~19.2 (P53) the
only new ingredient is the affine reflection $s_0$ in the wall
$\lambda+1 = k+2$; the finite reflection $s$ acts as before.
\end{solution}

\begin{solution}{66}
With $x = e^{i\phi}$, $\operatorname{ch}V_\lambda(x) = (x^{\lambda+1}-x^{-\lambda-1})/(x-x^{-1}) = \sin((\lambda+1)\phi)/\sin\phi$,
and $\phi = -\pi(\sigma+1)/(k+2)$ gives the ratio of sines (the sign of $\phi$
cancels). At $k = 2$, $K = 4$: $v_1 = (\sqrt2, 0, -\sqrt2)$, $v_2 = (1,-1,1)$,
$v_0 = (1,1,1)$ for $\sigma = 0,1,2$; then $v_1^2 = (2,0,2) = v_0+v_2$, the
Ising rule $1\times1 = 0+2$ evaluated pointwise. The evaluation matrix
$E_{\lambda\sigma} = S_{\lambda\sigma}/S_{0\sigma}$ is $S$ with its columns
rescaled by $1/S_{0\sigma}\ne0$, hence invertible since $S$ is unitary. By
Verlinde, $\sum_\nu N_{\lambda\mu}^{\ \nu}S_{\nu\sigma} = S_{\lambda\sigma}S_{\mu\sigma}/S_{0\sigma}$
(multiply the formula by $S_{\nu\sigma'}$ and sum over $\nu$ using unitarity),
so $\sum_\nu N_{\lambda\mu}^{\ \nu}v_\nu(\sigma) = v_\lambda(\sigma)v_\mu(\sigma)$:
fusion is pointwise multiplication of the vectors $v_\lambda$, and the
structure constants are recovered by changing back with $E^{-1}$.
\end{solution}

\begin{solution}{67}
(a) $Q^\vee = Q = \Z^r$ in root coordinates (simply laced), and $P = A^{-1}\Z^r$,
so $|P/KQ^\vee| = |A^{-1}\Z^r/K\Z^r| = |\Z^r/KA\Z^r| = K^r\det A$: $2K$ for $A_1$ and
$3K^2$ for $A_2$. (b) With $c = A^{-1}m_\mu/K$ ($m_\mu$ the Dynkin labels),
$\Theta_\mu(\tau,\xi) = \sum_{n\in\Z^r}f(n+c)$, $f(x) = e^{\pi iK\tau x^{\mathsf T}Ax+2\pi iKx^{\mathsf T}A\xi}$.
Poisson summation gives $\sum_nf(n+c) = \sum_{m\in\Z^r}\hat f(m)e^{2\pi im\cdot c}$
with
\begin{align*}
 \hat f(m) &= \int e^{\pi iK\tau x^{\mathsf T}Ax+2\pi ix^{\mathsf T}(KA\xi-m)}dx\\
 &= (-iK\tau)^{-r/2}(\det A)^{-1/2}\exp\!\big(-\tfrac{\pi i}{K\tau}(KA\xi-m)^{\mathsf T}A^{-1}(KA\xi-m)\big).
\end{align*}
Expanding, the exponent is $-\pi iK\xi^{\mathsf T}A\xi/\tau+2\pi im^{\mathsf T}\xi/\tau-\pi im^{\mathsf T}A^{-1}m/(K\tau)$,
and with $c' = A^{-1}m/K$ the last two terms are $2\pi iKc'^{\mathsf T}A(\xi/\tau)+\pi iK(-1/\tau)c'^{\mathsf T}Ac'$:
the summand of $\Theta_{\mu'}(-1/\tau,\xi/\tau)$ for the characteristic $m$.
Since $e^{2\pi im\cdot c} = e^{2\pi im^{\mathsf T}A^{-1}m_\mu/K}$ is constant on
cosets $m+KA\Z^r$, grouping gives
\[
 \Theta_\mu(\tau,\xi) = (-iK\tau)^{-r/2}(\det A)^{-1/2}e^{-\pi iK\xi^{\mathsf T}A\xi/\tau}\sum_{\mu'}e^{2\pi i\langle\mu,\mu'\rangle/K}\Theta_{\mu'}(-1/\tau,\xi/\tau).
\]
Replacing $(\tau,\xi)$ by $(-1/\tau,\xi/\tau)$, using $\Theta_{\mu'}(\tau,-\xi) = \Theta_{-\mu'}(\tau,\xi)$
and relabelling $\mu'\to-\mu'$ gives the stated formula, with
$(iK/\tau)^{-r/2}(\det A)^{-1/2} = (-i\tau)^{r/2}(K^r\det A)^{-1/2}$. For the
character ratio, numerator (level $k+h^\vee$) and denominator (level $h^\vee$)
share $(-i\tau)^{r/2}$, and their exponentials give
$e^{\pi i(k+h^\vee-h^\vee)\xi^{\mathsf T}A\xi/\tau} = e^{\pi ik\xi^{\mathsf T}A\xi/\tau}$;
the constants $|P/KQ^\vee|^{-1/2}$ are absorbed in the normalization of $S$.
Numerically, for $A_2$ at $K = 2$ the identity holds to $3\times10^{-15}$ over
the $12$ characteristics. (c) The lattice $\Z^r+\tau A\Z^r$ with the theta
functions above carries the polarization $KA$ (the quasi-periodicity factor
under $\xi\mapsto\xi+\tau e_j$ is $e^{-\pi iK\tau A_{jj}-2\pi iK(A\xi)_j}$), whose
section count is $\det(KA) = K^r\det A$, not the $K^r$ of a principal
polarization; so $\tau A$ is a period matrix for a nonprincipal bundle. The
genus-$r$ Siegel inversion $\Omega\mapsto-\Omega^{-1}$ would send $\tau A$ to
$-A^{-1}/\tau$, a different lattice, whereas the transformation derived above
stays on $\Omega = (-1/\tau)A$ and mixes the $|P/KQ^\vee|$ cosets.
\end{solution}

\begin{solution}{68}
Ng--Schauenburg: the torus representation factors through $\SL(2,\Z/N)$, a
finite group, for every modular tensor category, so no level and no algebra
gives infinitely many operators on the torus. Freedman--Larsen--Wang gives
density only for the braid group of a surface with enough labelled
punctures; universality therefore needs (i) punctures carrying alcove labels,
(ii) braid generators exchanging them. In the code, a puncture is a removed
top cell (the hole register: holes as removed top cells, holonomy read around
them), its label the sector of the harmonic space it carries, and a braid
generator a surgical move that exchanges two holes' positions; the
conformal block is the harmonic space of the punctured host at the register's
degree. None of these has been built as a labelled puncture with a fusion
space yet; the reading lists them as optional native constructions.
\end{solution}

\begin{solution}{69}
Hypothesis: the level-$k$ register of a boundary torus $M$ for a compact
simple $\mathfrak g$ is the theta space of the lattice $H^1(M;\Z)\otimes Q^\vee$
--- genus-$r$ theta functions with period matrix $\tau_M\cdot\operatorname{Gram}(Q^\vee)$,
the nonprincipal polarization of P67 and characteristics in
$P/KQ^\vee$, $K = k+h^\vee$ --- antisymmetrized
under $W$ and divided by $A_\rho$; the surface supplies only $H^1(M;\Z)$
and $\tau_M$, the algebra supplies $Q^\vee$, $W$ and $(\rho,h^\vee)$, and the
mapping class group of $M$ acts through $S$, $T$ on $\tau_M$ by the
Kac--Peterson matrices. The Weyl group acts on the fiber, not on the
triangulation: a $D_6$-symmetric torus realizes $D_6$ as mapping classes
of $M$ (elements of $\GL(2,\Z)$ acting on $H^1$), which coincides with
$W(G_2)$ as a group but is not the $W$ of the register, so its symmetry
buys nothing. First experiment (synthetic, the objective V5): at $\tau$ fixed,
build the antisymmetrized theta functions for $\mathfrak{su}(3)_k$ (and
$\mathfrak{su}(2)_k$ as the genus-one control), fit $S$ and $T$ from the
transformation law with the level-$k$ automorphy factor, and compare with
the Kac--Peterson closed forms and the Verlinde fusion; then take $\tau$ from
a seeded host and the flip-realized $S$, $T$ of the torus. Falsified if the
fit residual is not at machine precision, if the fitted matrices fail the
closed forms, or if the fusion coefficients are not the integers of
Kac--Walton. The geometric realization of the fiber would be a genus-$r$
boundary component on the locus $\Omega = \tau\cdot\operatorname{Gram}(Q^\vee)$.
\end{solution}

\part*{The experiment specification}

\section*{Whole harmonic states, pairings and elimination (Sections~26--28)}

\begin{solution}{70}
Necessity: if $B = TA$ then $\ker A\subseteq\ker B$, and ``defined on every
input'' means every input is $Ac$ for some $c$, i.e.\ $A$ is onto. Sufficiency:
let $C$ be a right inverse, $AC = I$. For any $c$, $A(c-CAc) = 0$, so
$c-CAc\in\ker A\subseteq\ker B$ and $Bc = BCAc$; hence $T := BC$ satisfies
$B = TA$. Two right inverses differ by a map $D$ with $AD = 0$, whose image
lies in $\ker A\subseteq\ker B$, so $BD = 0$ and $T$ is the same. Uniqueness:
$T$ is determined on $\operatorname{im}A$, which is everything. The
equivalent test is $B(I-CA) = 0$, since $I-CA$ projects onto $\ker A$.
Examples: $A = \begin{psmallmatrix}1&0\\0&0\end{psmallmatrix}$ (rank one) is
not onto: the input $e_2$ cannot be prescribed (incomplete coverage). With
$A = \begin{psmallmatrix}1&0\\0&0\end{psmallmatrix}$ and $B = \begin{psmallmatrix}0&0\\0&1\end{psmallmatrix}$,
$e_2\in\ker A$ but $Be_2 = e_2\ne0$: a nonzero output for zero input, so $T$
is ambiguous. The incomplete-coverage and ambiguous-output fixtures of NA2
exercise these, together with the hidden-mode fixture (an invisible class in
$\ker A\cap\ker B$, which is allowed) and the zero-map fixture.
\end{solution}

\begin{solution}{71}
With $\alpha = d_Uf$: $\alpha(v_1v_2) = g_{12}f(v_2)-f(v_1)$, $\alpha(v_0v_2) = g_{02}f(v_2)-f(v_0)$,
$\alpha(v_0v_1) = g_{01}f(v_1)-f(v_0)$, so
$(d_U\alpha)(v_0v_1v_2) = g_{01}g_{12}f(v_2)-g_{01}f(v_1)-g_{02}f(v_2)+f(v_0)+g_{01}f(v_1)-f(v_0)
= (g_{01}g_{12}-g_{02})f(v_2)$. This vanishes for all $f$ iff
$g_{01}g_{12}g_{02}^{-1} = 1$, the holonomy around the triangle. A gauge
transformation $f(v)\mapsto\lambda_vf(v)$, $g_{ij}\mapsto\lambda_ig_{ij}\lambda_j^{-1}$
preserves $d_U$; a chain coefficient $c(v_iv_j)$ must transform by
$\lambda_i^{-1}$ (the inverse system) for $\sum\alpha(e)c(e)$ to be invariant,
so the adjoint of $d_U$ is the boundary operator twisted by $U^{-1}$:
$\langle d_U\alpha,c\rangle = \langle\alpha,\partial^{U^{-1}}c\rangle$.
A cochain image $Z = G^U\Phi$ is closed iff it pairs to zero with every
$U^{-1}$-boundary, $(\partial_2^{U^{-1}})^{\mathsf T}Z = 0$; a chain frame
$\Phi$ is coclosed (a harmonic chain) iff its $U$-boundary vanishes,
$\partial_1^U\Phi = 0$, which is the adjoint statement $\delta_U\omega = 0$ in
chain coordinates. If $U$ is not flat, $d_U^2\ne0$ and there is no cohomology
to certify, which is why cohomological gluing needs a new justification there.
\end{solution}

\begin{solution}{72}
A boundary class with period vector $p$ is represented by the chain
$\Phi c$ with $Pc = p$, i.e.\ $c = P^{-1}p$; a dual class with periods $p^\vee$
by $\Phi^\vee(P^\vee)^{-1}p^\vee$. Their bilinear pairing is
\[
 (\Phi^\vee(P^\vee)^{-1}p^\vee)^{\mathsf T}M_{\rm own}\Phi P^{-1}p
= (p^\vee)^{\mathsf T}(P^\vee)^{-\mathsf T}B_0P^{-1}p,
\]
so $G_{\rm per} = (P^\vee)^{-\mathsf T}B_0P^{-1}$, and dropping $B_0$ is legitimate only when the frames are biorthonormal,
$B_0 = I$. Transport sends $p_A\mapsto p_B = Mp_A$ and $p^\vee_A\mapsto M^\vee p^\vee_A$;
the pairing is preserved iff $(M^\vee p^\vee)^{\mathsf T}G_B(Mp) = (p^\vee)^{\mathsf T}G_Ap$
for all $p,p^\vee$, i.e.\ $(M^\vee)^{\mathsf T}G_BM = G_A$. Under a gauge with
primal and dual factors $\rho,\tilde\rho$ the dressed mass transforms as
$G' = \tilde\rho^{-\mathsf T}G\rho^{-1}$, so $(\tilde\rho\tilde c)^{\mathsf T}G'(\rho c) = \tilde c^{\mathsf T}Gc$:
invariant. A primal--primal Gram transforms to $c^{\mathsf T}\rho^{\mathsf T}\tilde\rho^{-\mathsf T}G\rho^{-1}\rho c$,
which is not $c^{\mathsf T}Gc$ unless $\tilde\rho = \rho$, i.e.\ at zero
connection (P15).
\end{solution}

\begin{solution}{73}
The second block row of $Kx = 0$ is $K_{zb}b+K_{zz}z = 0$, so
$z = -K_{zz}^{-1}K_{zb}b$; substituting in the first row gives
$(K_{bb}-K_{bz}K_{zz}^{-1}K_{zb})b = 0$. Every solution is recovered from
$b$ by the formula for $z$, so nothing is discarded: the internal
coordinates are determined, not summed over. For the quadratic
$E(u) = \|u\|^2+\|r-T_2u\|^2$, $\nabla_{\bar u}E = u-T_2^\dagger(r-T_2u) = 0$
gives $(I+T_2^\dagger T_2)u = T_2^\dagger r$. Then $r-T_2u = (I-T_2(I+T_2^\dagger T_2)^{-1}T_2^\dagger)r = (I+T_2T_2^\dagger)^{-1}r$
(push-through identity), and
$E = r^\dagger[T_2(I+T_2^\dagger T_2)^{-2}T_2^\dagger+(I+T_2T_2^\dagger)^{-2}]r$;
writing $T_2(I+T_2^\dagger T_2)^{-2}T_2^\dagger = (I+T_2T_2^\dagger)^{-2}T_2T_2^\dagger$
the bracket is $(I+T_2T_2^\dagger)^{-2}(T_2T_2^\dagger+I) = (I+T_2T_2^\dagger)^{-1}$.
The zero-energy set is $r = 0$, i.e.\ $z = T_2T_1x$: composition. A partial
trace would replace the internal block by a mixed marginal, losing the
correlations between $b$ and $z$; the Schur complement keeps $z$ as a
function of $b$, so the joint solution is reconstructible.
\end{solution}

\section*{Quantum norms, channels and acceptance (Sections~29--30)}

\begin{solution}{74}
$\widehat T^\dagger\widehat T = Q_{\rm in}^{-1/2}T^\dagger Q_{\rm out}TQ_{\rm in}^{-1/2}$,
which is $I$ iff $T^\dagger Q_{\rm out}T = Q_{\rm in}$. With
$|\widehat T\rangle\!\rangle = \sum_{ij}\widehat T_{ij}|i\rangle_{\rm out}|j\rangle_{\rm ref}$,
$\Tr_{\rm out}|\widehat T\rangle\!\rangle\langle\!\langle\widehat T| = \sum_{jj'}\big(\sum_i\widehat T_{ij}\overline{\widehat T_{ij'}}\big)|j\rangle\langle j'|$,
and $\sum_i\overline{\widehat T_{ij'}}\widehat T_{ij} = (\widehat T^\dagger\widehat T)_{j'j}$,
so the marginal is $(\widehat T^\dagger\widehat T)^{\mathsf T}$ divided by the
norm $\|\widehat T\|_F^2 = \Tr\widehat T^\dagger\widehat T$; for an isometry this
is $I/d_{\rm in}$. The Schmidt decomposition of a vector in
$\C^{d_{\rm out}}\otimes\C^{d_{\rm in}}$ is the singular value decomposition of
its coefficient matrix, so the Schmidt coefficients are $\sigma_i/\|\widehat T\|_F$.
\end{solution}

\begin{solution}{75}
Each $|K_a\rangle\!\rangle\langle\!\langle K_a|$ is positive semidefinite, hence so
is the sum; conversely a positive $J$ has a spectral decomposition into such
terms, which is the Choi criterion for complete positivity. By P74,
$\Tr_{\rm out}|K_a\rangle\!\rangle\langle\!\langle K_a| = (K_a^\dagger K_a)^{\mathsf T}$,
so $\Tr_{\rm out}J = (\sum_aK_a^\dagger K_a)^{\mathsf T}$, which is $I$ iff the
Kraus operators are trace preserving. $V^\dagger V = \sum_{ab}K_a^\dagger K_b\langle a|b\rangle = \sum_aK_a^\dagger K_a = I$,
and $\Tr_E(V\rho V^\dagger) = \sum_{ab}K_a\rho K_b^\dagger\Tr(|a\rangle\langle b|) = \sum_aK_a\rho K_a^\dagger$.
Dimension count: $\dim(B\oplus E) = d_B+d_E$ while $\dim(B\otimes E) = d_Bd_E$;
a relation into a direct sum assigns to each input a pair (component in $B$,
component in $E$), and ``discarding'' $E$ is a projection onto a summand,
which is linear on vectors, whereas the partial trace over a tensor factor is
a map on operators that produces mixed states. No linear projection of a
direct sum reproduces $\Tr_E$, so a third boundary is not an environment in
the Stinespring sense.
\end{solution}

\begin{solution}{76}
(a) $H^2(T^2\times T^2) = H^0\otimes H^2\oplus H^1\otimes H^1\oplus H^2\otimes H^0$
has dimension $1+4+1 = 6$; the two degree-one registers occupy the
multidegree-$(1,1)$ summand $H^1\otimes H^1$, of dimension $4$, not all of $H^2$.
(b) $k^g$; for two tori, $k\cdot k = k^2$, the tensor product $\Theta_k\otimes\Theta_k$
(Proposition~11.1). (c) $V_1\otimes V_1 = V_0\oplus V_2$ has dimension $4$ and
two summands; at level $1$ the alcove is $\{0,1\}$ and the fusion is
$1\times1 = 0$ only, the summand $V_2$ being truncated by the affine wall
(P65). (d) A two-qubit gate has $4\times4 = 16$ amplitudes and its operator
vector lives in $\Theta_2\otimes\Theta_2\otimes\overline{\Theta_2\otimes\Theta_2}$,
also $16$; a direct sum of two two-dimensional port frames has $2+2 = 4$
amplitudes and its operators are block $2\times2$ maps with $4+4 = 8$ entries;
a selected-pair $\chi$ is one $2\times2$ matrix, $4$ entries. The state space
$\Theta_2\otimes\Theta_2$ also has $4$ amplitudes, but a tensor product is
more than a count: it comes with a factorization of the basis and of the
operator algebra ($D_{(p_1,p_2),(q_1,q_2)} = D_{p_1,q_1}\otimes D_{p_2,q_2}$), and
a direct sum has neither.
\end{solution}

\begin{solution}{77}
Oracle comparisons: R1--R4 (the read $M$ against the predicted
$\varphi^*$: identity or swap), C1 (the group law $\sigma^2 = I$), D1 (the
predicted direct sum), T4 (fitted $\rho_2(S),\rho_2(T)$ against the closed
forms $H$, $\operatorname{diag}(1,i)$), T5 and T6 (fitted matrices against the
Kronecker product and against controlled-$Z$), T7 (fits at two moduli
against each other), T8 (the geometric monodromy's fit against the
generator formulas), F2 (measured passes against $T_A,T_B,S$), F3 (the far
modulus against the bridge formula), F4 and T9 (the two orders against the
matrix products), G2 (the tube monodromy against the block form), E2 (the
decagon monodromy against the prediction from the side classes), E4--E5
(unitarity, order and the entangling witness against Clifford
enumeration), E6 (the product-collar control). Certificates: R5 (isotropy
for the declared form, with the flipped control as a negative), H1--H3
(closedness, coclosedness, independent span), J1--J3 (metric independence
and exactness of the difference), U1 (the transport defect reported, not
predicted), T1--T3 (fit residual, unitarity, projective group law on held-out
samples), F1 (the flipped host is a manifold with the declared boundary),
G1 and G3 (Betti numbers, isotropy, symmetry and positivity of the read
period matrix), G4--G5 (entropy of the geometric state and the degeneration
trend, reported), E1 and E3 (Euler characteristic, automorphism, fixed
period matrix).
\end{solution}

\section*{Neck geometry and encoded-state entropy (Section~33)}

\begin{solution}{78}
(a) The binary entropy is $h(p) = -p\log_2p-(1-p)\log_2(1-p)$ and
$-(1-p)\log_2(1-p) = (1-p)(p+p^2/2+\dots)/\ln2 = p/\ln2+O(p^2)$, so
$h(p) = -p\log_2p+p/\ln2+O(p^2)$. With $p = C\epsilon^2+O(\epsilon^3)$,
$-p\log_2p = -C\epsilon^2\log_2(C\epsilon^2)+O(\epsilon^3\log\epsilon)$, giving
the stated expansion; the $\log$ term dominates $C\epsilon^2$, so $S/\epsilon^2$
diverges as $\epsilon\to0$. (b) The exponent contains
$\pi ik(n+j/k)^{\mathsf T}\Omega(n+j/k)$ whose $\epsilon$-derivative is
$2\pi ik(n_1+j_1/k)(n_2+j_2/k)$; hence
$\partial_\epsilon\theta_j|_0 = \sum_{n_1,n_2}2\pi ik\,m_1m_2e^{(\cdots)_1}e^{(\cdots)_2}
= \frac{1}{2\pi ik}\Big(\sum_{n_1}2\pi ikm_1e^{(\cdots)_1}\Big)\Big(\sum_{n_2}2\pi ikm_2e^{(\cdots)_2}\Big)
= \frac{\theta'_{j_1}(z_1)\theta'_{j_2}(z_2)}{2\pi ik}$, since $\theta'_j(z) = \sum_n2\pi ik(n+j/k)e^{(\cdots)}$.
At level two, $\theta_0(-z) = \theta_0(z)$ (substitute $n\mapsto-n$) and
$\theta_1(-z) = \theta_1(z)$ (substitute $n\mapsto-n-1$: $-(n+\tfrac12)\mapsto n+\tfrac12$),
so $\theta'_j(0) = 0$ and the first-order change of the value column at
$z = 0$ vanishes. The normalized coherent state therefore has no
$O(\epsilon)$ component normal to the product manifold; the leading
entangling component is $O(\epsilon^2)$, the small singular value is
$O(\epsilon^2)$ and $p\sim C\epsilon^4$. (c) Build $\Omega(\epsilon)$, evaluate
\texttt{coherent\_state([0,0], Omega)} and
\texttt{entanglement\_entropy(vector, (2,2))} at $\epsilon = 0.04,0.02,0.01$;
the smaller probability falls by $16$ per halving (measured slope $3.999$ on a
log--log plot, values $4.5\times10^{-7}$, $2.8\times10^{-8}$, $1.8\times10^{-9}$),
while a random product state perturbed to first order gives slope $2$ and
the remainder $(S-[-p\log_2p+p/\ln2])/\epsilon^2\to0$.
\end{solution}

\part*{The measured constructions}

\section*{Reading the records (Section~35)}

\begin{solution}{79}
Expected counts: $17/17$ for each two-torus fixture, $27/27$ for the
four-torus fixtures; theta $8/8$ (two tori) and $10/10$ (four tori). The map:
R1--R4, rank and integrality of $M$ against $\varphi^*$: Theorem~4.2 and
$M = \varphi^*$ (Section~4); R5, isotropy for $\Omega_\partial$ with the flipped
control: Corollary~3.2 and Remark~3.3; H1--H3: the certificates that replace
Theorem~2.2 on complex masses; J1--J3: Theorem~2.4 (P3); C1:
Theorem~6.1's group law (P14); D1: Proposition~16.1; U1: Remark~7.1 and
Proposition~7.2; T1: Theorem~13.2 (closure of the theta space); T2:
Proposition~13.4; T3: the projective group law; T4: the $S$ and $T$ formulas
of Section~13; T5: Propositions~11.1 and~13.3; T6: the shear/controlled-$Z$
formula; T7: $\Omega$-independence of $\rho_k$; T8: the read monodromy acting
on the register. With an arbitrary closed two-dimensional band the periods
would still be periods of classes, so R1--R4, R5, C1, D1 and the theta checks
(which use only $M$) would pass; H2 (coclosedness) and H3 (independent kernel
of the same span) would fail, and J2--J3 would change meaning (the
representative would not be selected by the metric). This is why the harmonic
identification is certified separately from the topological reads.
\end{solution}

\begin{solution}{80}
(a) $L_b$ maps $(1,0)\mapsto(1,1)$, $(0,1)\mapsto(0,1)$, $(1,1)\mapsto(1,2)$;
the flip removes the row edges $(1,0)$ and adds the edges $(0,-1)$--$(1,1)$ of
displacement $(1,2)$, so the flipped edge set $\{(1,1),(0,1),(1,2)\}$ is exactly
$L_b$ of the grid's edge set: the flipped grid is the grid in the sheared
basis, relabelled by $(x,y)\mapsto(x,y+x)$. A cycle with new-coordinate
homology class $v$ has old class $L_bv$, so the far periods are read through
$L_b$ and the collar's relation is $M = L_b^{\mathsf T} = \begin{psmallmatrix}1&1\\0&1\end{psmallmatrix} = T_B$.
Likewise $L_a^{\mathsf T} = \begin{psmallmatrix}1&0\\1&1\end{psmallmatrix} = T_A$ and
$L_s^{\mathsf T} = \begin{psmallmatrix}0&-1\\1&0\end{psmallmatrix} = S$; with a
collar twist the passes compose as $M = L_{\rm last}^{\mathsf T}\cdots L_{\rm first}^{\mathsf T}\cdot M_{\rm twist}$.
(b) The second \texttt{b} pass flips the new-coordinate row edges, whose old
displacement is $L_b(1,0)^{\mathsf T} = (1,1)^{\mathsf T}$, and creates edges of old
displacement $L_b(1,2)^{\mathsf T} = (1,3)^{\mathsf T}$. On the $3$-periodic grid
$(1,3)\equiv(1,0)$: the new edge joins the same vertex pair as an original row
edge, which is still an edge of the host (it is the bottom of a first-pass
tetrahedron). Two edges on one vertex pair are not simplicial, so the code
refuses the pass by name. On the $5\times5$ grid $(1,3)\not\equiv(1,0)$, and two
passes fit.
(c) Expected: $20/20$ for one pass on grid $3$, $21/21$ for two passes on
grid $5$; F2 reads $T_B$, $T_A$, $S$ to $4\times10^{-15}$; F3 gives
$\tau' = (M_{21}+M_{22}\tau)/(M_{11}+M_{12}\tau)$ to $8\times10^{-16}$; F4 reads
$M = L_b^{\mathsf T}L_a^{\mathsf T} = T_BT_A = \begin{psmallmatrix}2&1\\1&1\end{psmallmatrix}$
for \texttt{ab} (pass \texttt{a} first, then \texttt{b}; the last pass acts
on the left) and $T_AT_B = \begin{psmallmatrix}1&1\\1&2\end{psmallmatrix}$ for
\texttt{ba}, which differ.
\end{solution}

\begin{solution}{81}
(a) Attaching a tube (an annulus $S^1\times[0,1]$, $\chi = 0$) to two tori
after removing a disk from each ($\chi = -1$ each) gives $\chi = -1-1+0 = -2$,
a connected closed orientable surface of genus two. $\chi(\partial W) = \chi(T^2)+\chi(T^2)+\chi(\Sigma_2) = -2$,
so $\chi(W) = -1$ (a compact odd-dimensional manifold has half the Euler
characteristic of its boundary). By P44 the tube is a $1$-handle joining
two collars, so $H^1(W) = H^1(W_1)\oplus H^1(W_2) = \Z^4$ and $b_1 = 4$; then
$1-4+b_2-b_3 = -1$ with $b_3 = 0$ (nonempty boundary) gives $b_2 = 2$:
$[1,4,2,0]$. $b_1(\partial W) = 2+2+4 = 8$ and the restriction has rank $4$
(Corollary~3.2). (b) The relation is a property of $H^1(W)$, which splits by
the $1$-handle argument: each harmonic class restricts to periods on one
collar's tori only, so $M = M_1\oplus M_2$ (measured $I_4$ or swap$\oplus$swap,
off-block norm $2.6\times10^{-15}$). The period matrix of $\Sigma_2$ is
defined by its own harmonic forms in the boundary metric, whose supports
cross the neck; $\Omega_{12}$ measures that leakage and is a metric quantity
of one connected surface, not of the relation. (c) Measured on the default
neck: $\operatorname{Im}\Omega$ positive with eigenvalues $0.83$, $1.08$,
$\Omega_{12} = 1.45\times10^{-4}-1.88\times10^{-4}i$, symmetric to $2.4\times10^{-15}$;
entropy $0$ for the diagonal control and $4\times10^{-7}$ bits mean; over the
$25$ necks $|\Omega_{12}|$ runs from $5\times10^{-8}$ to $2.4\times10^{-2}$ and
falls approximately exponentially with tube length, while the diagonal
approaches the two tori's moduli as the neck pinches. In the plumbing
description a genus-two surface near the separating degeneration is two tori
joined by an annulus of modulus $\ell/w$ (length over circumference), and
its period matrix is $\operatorname{diag}(\tau_1,\tau_2)$ plus corrections
linear in the plumbing parameter $t$, with $\Omega_{12}\propto t$ and
$|t|\sim e^{-2\pi\ell/w}$ (Fay's degeneration formula): exponentially small
in the neck's length-to-waist ratio, as observed.
\end{solution}

\begin{solution}{82}
The decagon with opposite sides identified ($E_{i+5}\sim E_i^{-1}$) is a
genus-two surface; coning the ten sides to a centre and barycentrically
subdividing gives $28$ vertices ($3$ identified corner/centre points, $5$ side
midpoints, $10$ spoke midpoints, $10$ barycentres), $90$ edges and $60$
faces, $28-90+60 = -2$. The rotation permutes the vertices and preserves the
face list, so \texttt{is\_automorphism} returns true; it has order ten as a
simplicial isometry of the flat decagon metric. \texttt{induced\_map(2)}
returns $M$ on $(A_1,A_2,B_1,B_2)$ with $\det M = 1$, $M^{\mathsf T}JM = J$ and
$M^5 = I$ (P47), mixing the two factors. The fit
\texttt{fit = ThetaRegister(level=2).weil(M, omega)} returns
\texttt{fit.matrix} unitary with residual $\sim6\times10^{-16}$ and
$\rho^5\propto I$ to $1.5\times10^{-15}$;
\texttt{ThetaRegister.schmidt\_rank(fit.matrix, (2,2))} is $4$. Applying
\texttt{fit.apply} to the $36$ products of $X,Y,Z$ eigenstates and calling
\texttt{entanglement\_entropy(v, (2,2))} gives exactly one bit for $16$ of
them and zero for the rest; the gate is at projective distance $1.00$ from
every $A\otimes B$ and $0.77$ from every $\mathrm{SWAP}\cdot A\otimes B$ over
the $24\times24$ single-qubit Cliffords (E5). The CLI run reports $13/13$;
$r^4$ gives order five and Schmidt rank two, $r^5 = -I$ Schmidt rank one and
no entangling power, and the product collar reads $M = I$ with entropy
$5\times10^{-29}$ (E6).
\end{solution}

\begin{solution}{83}
(a) With the flat metric of modulus $\tau$ the cotangent weights are those of
the flat torus; the harmonic cochains are the constant forms $dx$, $dy$, and
the holomorphic combination $dx+\tau\,dy$ (P20 and Section~26) has
$A$-period $1$ and $B$-period $\tau$, so $\Omega = \tau$; the cup pairing
$C_{ab} = \sum_{\rm faces}\operatorname{area}\,(W_a\times W_b)$ evaluated on the
harmonic basis gives the intersection matrix $-PC^{-1}P^{\mathsf T} = J$ with
$A\cdot B = +1$, the sign the code calibrates on this case. (b)
\texttt{s.periods.omega} is symmetric to machine precision with
$\operatorname{Im}\Omega$ positive definite (eigenvalues $0.46$, $1.66$); with
$P,Q,R,S$ the blocks of \texttt{s.induced\_map(2)}, $(R+S\Omega)(P+Q\Omega)^{-1}$
equals $\Omega$ to $1.5\times10^{-15}$ (E3). A simplicial isometry
$\varphi$ of the surface preserves the metric, hence the Hodge star and the
space of holomorphic forms; it acts on $H^1(\Sigma;\Z)$ by $M_{\rm per}$ and
on the normalized holomorphic basis by the change of marking of P36,
so the period matrix in the transported marking is $\Omega' = \Omega$: $\Omega$
is a fixed point of $\mathsf M$ on the Siegel upper half-space, which is why
the read matrix is fixed by $M$. Being fixed by an order-five element
singles out a special point of $\mathcal H_2$ (the curve $y^2 = x^5-1$), but
nothing forces it onto the locus $\Omega = \tau A$ with $A$ a root Gram
matrix; that locus is a different one-parameter family, and the root register
also needs the nonprincipal polarization $KA$ (P67), which a
principally polarized Jacobian does not carry.
\end{solution}

\end{document}
